%% parts/p03/classification-theory.tex --- P3 prime/atomic models,
%% indiscernibles, and the Phase-4 classification-theory survey.
\chapter{Prime Models and Classification Theory}
\label{sec:p03-classification-theory}

\noindent\textit{Reader note.}  This is a separately authored survey after the
exact inherited model-theory chapters.  Its prime/atomic, indiscernibility, and
classification entries follow \cite{Marker2002,Hodges1993}; they do not replace
or alter the inherited categoricity theorems.

This chapter fixes the minimal vocabulary needed to distinguish prime and
atomic models and the principal modern dividing lines.  The survey entries do
not replace the hypotheses of the theorems they cite; in particular,
completeness, countability, parameter sets, and the chosen ambient model are
always explicit.

Throughout this chapter,
\[
  \abs T\defeq\max(\abs\Lang,\aleph_0)
\]
for a complete \(\Lang\)-theory \(T\).  This is a metalinguistic size
abbreviation, not the cardinality of the set of theorems after an arbitrary
coding; it is used in the EM spectrum and local-character bounds below.

\section{Fra\"iss\'e classes and the random graph}
\label{sec:p03-fraisse}

\begin{hbdefinition}{def:p03-age-fraisse}{Ages and Fra\"iss\'e classes}
Fix a countable relational signature \(\tau\).  The \emph{age}
\(\operatorname{Age}(M)\) of a countable \(\tau\)-structure \(M\) is the
class, up to isomorphism, of its finite induced substructures.  A nonempty
isomorphism-closed class \(\mathcal K\) of finite \(\tau\)-structures is a
\emph{Fra\"iss\'e class} when it has countably many isomorphism types,
contains structures of arbitrarily large finite cardinality, and satisfies:
\begin{enumerate}
  \item the \emph{hereditary property}: every induced substructure of a
        member of \(\mathcal K\) belongs to \(\mathcal K\);
  \item the \emph{joint-embedding property}: any two members embed into one
        member of \(\mathcal K\); and
  \item the \emph{amalgamation property}: for embeddings
        \(f_i:A\to B_i\) with \(A,B_1,B_2\in\mathcal K\), there are
        \(C\in\mathcal K\) and embeddings \(g_i:B_i\to C\) satisfying
        \(g_1f_1=g_2f_2\).
\end{enumerate}
The age is needed to record exactly the finite configurations of one
structure; a Fra\"iss\'e class is instead a class satisfying the three
extension conditions from which a structure can be reconstructed.  Both are
metalinguistic classes of finite structures, not first-order theories.  Their
difference is used by the existence-and-uniqueness theorem below.
\index{age of a structure}\index{Fraisse class@Fra\"iss\'e class}
\index{amalgamation property}
\end{hbdefinition}

\begin{hbdefinition}{def:p03-ultrahomogeneous}{Ultrahomogeneous structure}
A countable relational structure \(M\) is \emph{ultrahomogeneous} when every
isomorphism between finite induced substructures of \(M\) extends to an
automorphism of \(M\).  This is an extension property for finite
\emph{isomorphisms}; it is stronger than the order-indiscernibility of a
chosen sequence and different from the elementary-map homogeneity of
Definition~\ref{def:homogeneity}.  The distinction is needed because the
Fra\"iss\'e limit is characterized by this finite structural extension
property.
\index{ultrahomogeneous structure}
\end{hbdefinition}

\begin{hbtheorem}{thm:p03-fraisse-limit}{cited=Fraisse1954}{Fra\"iss\'e's theorem}
Every Fra\"iss\'e class \(\mathcal K\) is the age of a countably infinite
ultrahomogeneous structure \(F\).  Any two countable ultrahomogeneous
structures having age \(\mathcal K\) are isomorphic.
\end{hbtheorem}

For existence, enumerate all finite embedding problems in \(\mathcal K\) and
build an increasing chain that solves the next problem at each stage; JEP
starts the chain and AP makes each extension compatible with what has already
been built.  Its union has age \(\mathcal K\) and the required one-point
extension property.  For uniqueness, enumerate two candidate limits and run
back-and-forth, using that extension property to enlarge each finite partial
isomorphism.  This also explains why HP, JEP, and AP do different work in the
theorem rather than serving as three names for one closure condition
\cite{Fraisse1954,Hodges1993}.

\begin{hbtheorem}{thm:p03-random-graph}{cited=Rado1964}{The countable random graph}
There is, up to isomorphism, exactly one countable simple graph \(R\) with the
following extension property: whenever \(U,V\subseteq R\) are disjoint and
finite, some vertex outside \(U\cup V\) is adjacent to every vertex in \(U\)
and to no vertex in \(V\).  It is the Fra\"iss\'e limit of the class of all
finite simple graphs, is ultrahomogeneous, and its complete theory is
\(\aleph_0\)-categorical.
\end{hbtheorem}

Indeed, all finite graphs have HP, JEP by disjoint union, and AP by taking the
free amalgam with no new cross-edges.  The displayed extension property is
the resulting one-point extension property.  Back-and-forth proves uniqueness
of its countable model; equivalently, the extension axioms together with the
graph axioms form an \(\aleph_0\)-categorical theory.  This deterministic
characterization is what ``the random graph'' means below; it does not name a
probability distribution on finite graphs \cite{Rado1964,Marker2002}.
\index{random graph}

\section{Prime and atomic models}
\label{sec:p03-prime-atomic}

\begin{hbdefinition}{def:p03-atomic-prime}{Atomic and prime models}
Let \(T\) be a complete satisfiable theory.  A model \(\mathcal M\entails T\) is
\emph{atomic} when every finite tuple from \(\dom\mathcal M\) realizes an
isolated complete type over the empty set, where ``isolated'' is the
principal-complete-type condition of
Definition~\ref{def:principal-type}.  It is \emph{prime} when, for every
\(\mathcal N\entails T\), there is an elementary embedding
\(\mathcal M\to\mathcal N\).  Atomicity classifies which types are realized
inside one model; primeness compares that model with every model of \(T\).
The notions are therefore not definitionally redundant, and the countability
theorem below is what identifies them in its stated range.
\index{atomic model}
\index{prime model}
\end{hbdefinition}

\begin{hbtheorem}{thm:p03-prime-atomic}{proved=annex:p3.vaught.overview-prime}{Prime--atomic model package}
Let \(T\) be a complete satisfiable theory in a countable language.
\begin{enumerate}
  \item For every \emph{countable} \(\mathcal M\entails T\), the model
        \(\mathcal M\) is prime if and only if it is atomic.
  \item The theory \(T\) has a prime model if and only if, for every
        \(n<\omega\), the isolated types are dense in \(S_n(T)\).
\end{enumerate}
The countability restriction belongs to clause~1: an arbitrary atomic model
need not be prime.
\end{hbtheorem}

The complete construction below supplies the prime-to-atomic direction by
omitting a nonisolated type, while the atomic-to-prime direction enumerates
the countable model and realizes successive isolating formulas in the target.
The overview proof treats finite models and function languages separately;
\cite{Marker2002} gives its independently attributed presentation.
The Ryll--Nardzewski theorem stated in
Section~\ref{sec:qe-categoricity-02} gives the sharper package: if the
countable complete theory \(T\) has infinite
models, then \(T\) is \(\aleph_0\)-categorical exactly when every
\(S_n(T)\) is finite, equivalently when every complete \(n\)-type is isolated.
In that setting the unique countable model is both prime and saturated
\cite{RyllNardzewski1959,Engeler1959,Svenonius1959}.  None of these conclusions
is asserted after dropping completeness or the displayed countability
hypotheses.

\section{Vaught's Formula-Algebra Construction of Countable Models}
\label{sec:p03-n-a-003-01}

Vaught's original account begins with Boolean algebras of formulas and
their dual prime ideals. It constructs models that simultaneously omit
specified types and control which formulas their finite tuples satisfy.
The elementary embedding arguments then identify prime and atomic models,
and the saturation construction treats a countable family of relative
interpretations at once \cite[Secs.~1--4, pp.~304--314]{Vaught1961}.
This presentation keeps those constructions and their original
countability conditions explicit.

\subsection{The source's language and formula algebras}
\label{sec:p03-vaught-formula-algebras}

\begin{hbconvention}{def:p03-vaught-standing}{Standing theory and meaning of countability}
Unless a result explicitly changes the setting, \(T\) is a complete
satisfiable theory in an at most countable language \(L\) of finite-arity
relations and individual constants, with logical identity, and has an
infinite model. Completeness then implies that every model is infinite.
Vaught's denumerable models are countably infinite.
The variables are \(v_0,v_1,\ldots\); all substitutions avoid capture.
The constructions use the ambient Choice assumption of
\cref{thm:fol-compactness}.
For theories in these constructions, consistency is satisfiability;
first-order completeness supplies the equivalent syntactic condition.
Countable enumeration here is a set-theoretic enumeration; no algorithm
for deciding consistency or finding isolators is asserted.
The source prints \(\bigvee v_i\) for existential and
\(\bigwedge v_i\) for universal quantification; the displays here
write \(\exists v_i\) and \(\forall v_i\), with the same binding
and truth clauses. These quantifier signs do not denote infinitary
Boolean disjunction or conjunction.
\end{hbconvention}

\begin{hbdefinition}{def:p03-vaught-formula-algebra}{Formulas, validity, atoms and atomistic algebras}
Let \(F_n(T)\) be the formulas with free variables among
\(v_0,\ldots,v_{n-1}\). In this source presentation
\(\vdash_T\varphi\) means validity of the universal closure in every
model of \(T\). Write \(\varphi\approx_T\psi\) when
\(\vdash_T(\varphi\leftrightarrow\psi)\).
The quotient \(\mathbb B_n(T)=F_n(T)/{\approx_T}\), with conjunction,
disjunction and negation, is a Boolean algebra.
We also use the source's \(F_n(T)\) for the corresponding formula
prealgebra, retaining its equivalence relation.

A formula is nonzero, or consistent with \(T\), when
\(T\cup\{\exists v_0\cdots\exists v_{n-1}\varphi\}\) is satisfiable.
An \emph{atom} is a consistent \(\alpha\in F_n(T)\) such that,
for each \(\varphi\in F_n(T)\), consistency of
\(\alpha\land\varphi\) implies \(\vdash_T\alpha\to\varphi\).
An \emph{atomless element} is a consistent \(\beta\) with no atom
\(\alpha\) satisfying \(\vdash_T\alpha\to\beta\).
The algebra is \emph{atomistic} if every consistent formula has an
atom below it. These properties refer to the quotient order; they
do not say that a particular tuple realizes a type.
\end{hbdefinition}

\begin{hbdefinition}{def:p03-vaught-prime-ideals}{The source's dual prime ideals and realized tuples}
A dual prime ideal \(P\) is a nonempty proper subset of \(F_n(T)\),
closed under conjunction and upward under \(T\)-implication, that
contains either \(\varphi\) or \(\neg\varphi\) for every \(\varphi\).
Exactly one of those alternatives belongs to \(P\).
The source calls this a prime ideal; its direction is the one usually
called an ultrafilter. Let \(\mathcal P_n(T)\) be the set of these \(P\).
It is principal if
\[
 P=\{\varphi:\vdash_T\alpha\to\varphi\}
\]
for some consistent \(\alpha\). For \(\mathcal M\models T\), put
\[
 P(\mathcal M)=\{\bar a\in(\dom\mathcal M)^n:
                         \mathcal M\models\varphi(\bar a)
                         \text{ for every }\varphi\in P\}.
\]
Repeated coordinates are allowed.

With distinct new constants \(c_0,\ldots,c_{n-1}\), write
\(\overline\varphi=\varphi(c_0,\ldots,c_{n-1})\) and
\(\overline P=T\cup\{\overline\varphi:\varphi\in P\}\).
The latter is a complete theory in the language with those named
coordinates.
\end{hbdefinition}

\begin{hblemma}{thm:p03-vaught-algebra-types}{proved=annex:p3.vaught.algebra-types}{Atoms, complete types and named expansions}
The dual prime ideals are exactly the consistent complete \(n\)-types
over the empty set, identifying \(\mathcal P_n(T)\) with \(S_n(T)\).
Their source principality agrees with \cref{def:principal-type}.
A principal complete type contains an atom and is isolated by it.
For a \emph{consistent} \(\alpha\), being an atom is equivalent to all
its realizations yielding elementarily equivalent named expansions.
The basic sets \([\varphi]=\{P:\varphi\in P\}\) form the Stone topology;
atomisticity is equivalent to density of isolated types.
\(F_0(T)\) has exactly two equivalence classes and one complete type.
\end{hblemma}

\subsection{The simultaneous Henkin and omission construction}
\label{sec:p03-vaught-simultaneous}

\begin{hbdefinition}{def:p03-vaught-henkin-data}{New constants and witness requirements}
Let \(C=\{c_i:i<\omega\}\) be distinct symbols outside \(L\).
At a finite stage let \(\Gamma_s\) be the conjunction of all accumulated
sentences, consistent with the starting theory.
Enumerate every formula \(\varphi_s(v_0)\) of \(L(C)\) with at most
the displayed free variable. Choose \(c_{\nu_s}\) with least index
absent from \(\varphi_s,\Gamma_s\), and use
\[
 W_s=(\exists v_0\varphi_s)\to\varphi_s(c_{\nu_s}).
\]
Different constants may denote the same element.
An existential block of length zero makes no quantification.
\end{hbdefinition}

\begin{hblemma}{thm:p03-vaught-henkin-naming}{proved=annex:p3.vaught.henkin-naming}{Fresh witnesses, projections and a fully named model}
Adding \(W_s\) preserves consistency at the finite stage.
Replacing distinct new constants by free variables, after freshening
binders, equates consistency of their sentence instance with consistency
of its existential closure.
Any consistent theory in \(L(C)\) containing witnesses for every such
formula has a model whose \(L\)-domain consists of the \(C\)-values.
The named part is elementary in a model of that witness theory.
Thus every consistent countable relational/constant theory has an
at most countable model, and an infinite one when its axioms force
infinitude.
In the fresh-witness construction each \(C\)-value has infinitely
many distinct constant symbols naming it.
\end{hblemma}

\begin{hbtheorem}{thm:p03-vaught-simultaneous-construction}{proved=annex:p3.vaught.simultaneous}{The full original simultaneous construction}
Let \((P_j)_{j\in J}\), \(J\subseteq\omega\), be any specified finite
or countably infinite family of nonprincipal complete types
\(P_j\in\mathcal P_{r_j}(T)\), with finite arities.
The source's \(p_j+1\) is this total arity \(r_j\).
There is a countably infinite \(\mathcal M\models T\) such that:
\begin{enumerate}
\item Every finite tuple satisfies either an atom or an atomless formula
      of its corresponding \(F_n(T)\).
\item \(P_j(\mathcal M)=\varnothing\) for every \(j\in J\).
\end{enumerate}
Either condition may be imposed alone. Arity zero needs no omission
requirement, since its sole complete type is principal.
This is the simultaneous content of source Theorem~2.1, not merely
the existence of a model omitting one type.
\end{hbtheorem}

\begin{hbmethod}{def:p03-vaught-two-branch-construction}{The source's two branches at one stage}
Enumerate every requirement \(\Pi_s=(P,d_0,\ldots,d_{r-1})\), where the
\(d_i\) are distinct \(C\)-symbols and either \(P=\varnothing\) or
\(P\) is a listed type of arity \(r\geq1\).
Pair this enumeration with the witness enumeration.
Let \(\gamma_s=\Gamma_s\land W_s\).
Write \(\bar d\) for the task tuple and \(\bar e\) for the other
occurring \(C\)-symbols, disjoint from \(\bar d\).
Replace these two tuples by free slots \(\bar v,\bar w\), after
freshening binders, obtaining \(\theta_s(\bar v,\bar w)\) in \(L\).
Put \(\beta_s(\bar v)=\exists\bar w\,\theta_s(\bar v,\bar w)\).
It is consistent.

If an atom \(\alpha\) lies below \(\beta_s\), add
\(\alpha(\bar d)\land\gamma_s\).
If \(\beta_s\) is atomless and \(P=\varnothing\), add \(\gamma_s\).
If \(\beta_s\) is atomless and \(P\ne\varnothing\), choose
\(\delta\notin P\) with \(\beta_s\land\delta\) consistent and add
\(\delta(\bar d)\land W_s\), retaining all previous sentences.
The annex proves that this choice always exists and that each branch
preserves consistency. The infinitely many aliases, rather than a
restriction to tuples of distinct elements, allow these syntactically
distinct constant requirements to cover every tuple.
\end{hbmethod}

\begin{hbcorollary}{thm:p03-vaught-automatic-omission}{proved=annex:p3.vaught.automatic-omission}{Further omitted types in the closure of isolated points}
A model with the first property of the simultaneous theorem omits
every nonprincipal complete type whose every neighborhood contains
an isolated point. In this countable formula language these are exactly
the nonprincipal points which are limits of a sequence of distinct
isolated types. This automatically omitted collection may be uncountable;
the separately specified omission list remains countable.
\end{hbcorollary}

\subsection{Elementary maps, prime models and atomic models}
\label{sec:p03-vaught-prime-atomic}

\begin{hbdefinition}{def:p03-vaught-model-notions}{Atomicity, primeness, universality and homogeneity}
An \(\mathcal A\models T\) is atomic when each finite tuple satisfies
an atom of its formula algebra. It is prime when it elementarily embeds
in every model of \(T\). These agree with
\cref{def:p03-atomic-prime} by the algebra/type identification.
The elementary-embedding condition is Vaught's stated prime notion;
his introductory footnote~3 distinguishes earlier ordinary-embedding
variants.
It is \(\aleph_0\)-universal, in the source's convention, when it is
countably infinite and every countably infinite \(T\)-model
elementarily embeds in it.

It is homogeneous when any two finite tuples with the same complete
type are carried one to the other by an automorphism.
Here the data are full first-order types, not just finite relational
isomorphism patterns. An elementary tuple map preserves the truth of
every formula, including identity, with those tuple entries as parameters.
Repeated entries therefore must correspond to repeated entries.
\end{hbdefinition}

\begin{hblemma}{thm:p03-vaught-prefix-embedding}{proved=annex:p3.vaught.prefix-embedding}{The source's complete-prefix embedding criterion}
Suppose \(\dom\mathcal A=\{a_i:i<\omega\}\), allowing repetitions,
and \(\bar b=(b_i)_{i<\omega}\) is in \(\mathcal B\models T\).
If, for every \(n\), \((a_0,\ldots,a_{n-1})\) and
\((b_0,\ldots,b_{n-1})\) have the same complete type, then
\(a_i\mapsto b_i\) is well defined and injective, respects every
constant and relation, and is an elementary embedding.
\end{hblemma}

\begin{hblemma}{thm:p03-vaught-atomic-parameters}{proved=annex:p3.vaught.atomic-parameters}{Naming a finite tuple preserves atomicity}
If \(\mathcal A\) is atomic and \(\bar a\) is a finite tuple from it,
the expansion \((\mathcal A,\bar a)\), considered as a model of its
complete named theory, is atomic. Isolation in that expansion is
over the specified finite parameters; it is not isolation over
an arbitrary infinite parameter set.
\end{hblemma}

\begin{hblemma}{thm:p03-vaught-atomic-embedding}{proved=annex:p3.vaught.atomic-embedding}{Original Lemma 3.1 with its prescribed tuple}
Let \(\mathcal A\) be countably infinite and atomic, and let
\(\mathcal B\) be any \(T\)-model, of any cardinality.
There is an elementary embedding \(\mathcal A\to\mathcal B\).
If finite tuples \(\bar a,\bar b\) have the same complete type,
the embedding can be chosen to send \(\bar a\) to \(\bar b\).
The construction uses atoms of successive source tuples; it imposes
no atomicity or saturation hypothesis on the target.
\end{hblemma}

\begin{hbtheorem}{thm:p03-vaught-atomic-uniqueness}{proved=annex:p3.vaught.atomic-uniqueness}{Original Theorem 3.2: uniqueness by full back-and-forth}
Two countably infinite atomic models of \(T\) are isomorphic.
An elementary correspondence between prescribed finite tuples
can be extended to such an isomorphism.
\end{hbtheorem}

\begin{hbtheorem}{thm:p03-vaught-atomic-homogeneity}{proved=annex:p3.vaught.atomic-homogeneity}{Original Theorem 3.3: elementary homogeneity}
Every countably infinite atomic model of \(T\) is homogeneous
in the stated complete-type sense.
\end{hbtheorem}

\begin{hbtheorem}{thm:p03-vaught-prime-atomic}{proved=annex:p3.vaught.prime-atomic}{Original Theorem 3.4: prime exactly when countable and atomic}
A model of the standing \(T\) is prime if and only if it is
countably infinite and atomic.
In particular, atomicity of an uncountable model does not imply primeness.
Any two prime models are isomorphic.
\end{hbtheorem}

\begin{hbtheorem}{thm:p03-vaught-isolation-density}{proved=annex:p3.vaught.isolation-density}{Original Theorem 3.5: the full existence equivalence}
The following are equivalent:
\begin{enumerate}
\item \(T\) has a prime model.
\item \(T\) has an atomic model, without a cardinal restriction here.
\item Every \(F_n(T)\) is atomistic.
\item For every \(n<\omega\), isolated types are dense in \(S_n(T)\).
\end{enumerate}
The atomic model in clause 2 can have any size; the existence of a
countably infinite atomic model then follows from the simultaneous
construction. This existence equivalence does not assert that the
original arbitrary-sized atomic model itself is prime.
\end{hbtheorem}

\begin{hblemma}{thm:p03-vaught-nonprincipal-algebra}{proved=annex:p3.vaught.nonprincipal-algebra}{Original fact 3.6 with its complete algebra proof}
The countable Boolean algebra \(\mathbb B_n(T)\) has infinitely many
elements if and only if it has a nonprincipal dual prime ideal.
Equivalently, infinitely many inequivalent \(n\)-variable formulas
force a nonisolated complete \(n\)-type.
\end{hblemma}

\begin{hbtheorem}{thm:p03-vaught-finite-formulas-categoricity}{proved=annex:p3.vaught.finite-formulas-categoricity}{The source section 3 consequence of the construction}
For the standing complete countable-language \(T\), the following
are equivalent:
\begin{enumerate}
\item All countably infinite models of \(T\) are isomorphic.
\item Each \(F_n(T)\) has only finitely many inequivalent formulas.
\item Each \(\mathcal P_n(T)\) is finite.
\item Every complete finite-tuple type is principal.
\end{enumerate}
This is the finite-formula characterization discussed after source
Theorem~3.5; its proof uses the supplied atomic and omission constructions.
The original Ryll-Nardzewski, Engeler and Svenonius presentations
remain distinct from this derivation in Vaught's account.
\end{hbtheorem}

\subsection{The exact saturation and parameter scopes}
\label{sec:p03-vaught-saturation}

\begin{hbdefinition}{def:p03-vaught-saturation}{Weak saturation and the source's one-coordinate extension}
A model is weakly saturated when \(P(\mathcal M)\ne\varnothing\)
for every \(P\in\mathcal P_n(T)\) and every finite \(n\).
It is saturated in Vaught's sense when it is weakly saturated and,
whenever \(P\in\mathcal P_n(T)\), \(Q\in\mathcal P_{n+1}(T)\),
\(P\subseteq Q\) under the inclusion of the first \(n\) variable slots,
and \(\bar a\in P(\mathcal M)\), some \(b\) satisfies
\((\bar a,b)\in Q(\mathcal M)\).
Since \(P\) is complete, \(P\subseteq Q\) means
\(Q\cap F_n(T)=P\).

In modern parameter notation, \(\kappa\)-saturation means realization
of every complete finite-tuple type over each parameter set \(A\)
of cardinality less than the specified infinite \(\kappa\).
For \(\kappa=\aleph_0\), this concerns finite parameter sets.
Homogeneity remains the automorphism property of finite elementary
tuple maps defined above; it is not introduced as a synonym for saturation.
\end{hbdefinition}

\begin{hblemma}{thm:p03-vaught-saturation-parameters}{proved=annex:p3.vaught.saturation-parameters}{The finite-parameter meaning of the source definition}
Vaught saturation is equivalent to realization of every complete
finite-tuple type over every finite parameter set, relative to the
complete theory with those parameters named.
After naming an \(m\)-tuple, its complete \(r\)-types inject into
\(\mathcal P_{m+r}(T)\), extending that tuple's fixed type.
Thus smallness of the empty-parameter type spaces also bounds
the spaces over each finite parameter tuple.
For an infinite countable model this does not imply realization
over arbitrary countable parameter sets: a consistent type asking
for a new element different from every existing element is omitted.
\end{hblemma}

\subsection{Relative interpretations and the full original construction}
\label{sec:p03-vaught-relative-saturation}

\begin{hbdefinition}{def:p03-vaught-relative-interpretation}{The source's definable-domain interpretations}
Here \(T_1\) denotes a countable relational/constant target theory.
The interpretation formulas have no external parameters; fixed
parameters must first be named in that target language.
A relative interpretation \(I\) of \(T\) in \(T_1\) specifies a
formula \(D(x)\) for a nonempty domain, a formula \(\rho_R(\bar x)\)
for each relation of \(L\), and a unique distinguished element
in that domain for each individual constant.
Require \(T_1\models\exists x\,D(x)\).
The distinguished element may be given by a target constant or
by a formula \(K_d(x)\) with
\[
                 T_1\models\exists!x\,(D(x)\land K_d(x)).
\]
Equality is actual identity on the domain, not a quotient relation.
Translate quantifiers by
\[
 (\exists x\varphi)^I=\exists x(D(x)\land\varphi^I),\qquad
 (\forall x\varphi)^I=\forall x(D(x)\to\varphi^I),
\]
relations by their specified formulas, and Boolean operations
compositionally. For a defined distinguished element, use a fresh
existential variable satisfying its unique defining formula in the
atomic translation. This also treats repeated constant occurrences
coherently. The interpretation requirement is that
\(T_1\models\sigma^I\) for every axiom \(\sigma\) of \(T\).

In \(\mathcal B\models T_1\), denote the resulting \(T\)-structure
on \(D(\mathcal B)\) by \(I(\mathcal B)\).
These are the definable-domain relative interpretations of source
section 1, with distinguished-element details written explicitly.
\end{hbdefinition}

\begin{hblemma}{thm:p03-vaught-interpretation-preservation}{proved=annex:p3.vaught.interpretation-preservation}{Truth, elementary maps and unions under interpretation}
For assignments in the interpreted domain,
\[
 I(\mathcal B)\models\varphi(\bar a)
                 \quad\Longleftrightarrow\quad
 \mathcal B\models\varphi^I(\bar a).
\]
If \(\mathcal B\preccurlyeq\mathcal B'\), then
\(I(\mathcal B)\preccurlyeq I(\mathcal B')\).
For an elementary increasing chain, interpretation commutes
with its union, including the defined distinguished elements.
\end{hblemma}

\begin{hblemma}{thm:p03-vaught-interpreted-extension}{proved=annex:p3.vaught.interpreted-extension}{Original Lemma 4.1: one extension with its elementary diagram}
Let \(\mathcal B\models T_1\) be countably infinite,
\(\mathcal A=I(\mathcal B)\), and
\(\bar a\in P(\mathcal A)\) with
\(P\in\mathcal P_n(T)\), \(P\subseteq Q\in\mathcal P_{n+1}(T)\).
There is a countably infinite elementary extension
\(\mathcal B^*\succeq\mathcal B\) such that, in
\(\mathcal A^*=I(\mathcal B^*)\), some \(b\) has
\((\bar a,b)\in Q(\mathcal A^*)\).
The empty-prefix case \(n=0\) is included.
\end{hblemma}

\begin{hbtheorem}{thm:p03-vaught-saturated-existence}{proved=annex:p3.vaught.saturated-existence}{Full original Theorem 4.2 for a family of interpretations}
Suppose every \(\mathcal P_n(T)\) is countable.
Let \(T_1\) be a consistent countable theory and let
\((I_i)_{i\in J}\), for a nonempty finite or countable \(J\), be
relative interpretations of \(T\) in \(T_1\).
There is a countably infinite \(\mathcal B\models T_1\) for which
every \(I_i(\mathcal B)\) is a saturated \(T\)-model.
The construction first fills every extension requirement of one
countable stage, and then repeats this operation over an outer
elementary chain to cover new tuples.
The target theory need not be complete.
With the identity interpretation this supplies a countably infinite
saturated model of \(T\).
\end{hbtheorem}

\subsection{Universality, uniqueness and the smallness criterion}
\label{sec:p03-vaught-universality}

\begin{hblemma}{thm:p03-vaught-saturated-embedding}{proved=annex:p3.vaught.saturated-embedding}{Original Lemma 4.3 with a prescribed finite map}
Let \(\mathcal B\) be countably infinite and saturated, and let
\(\mathcal A\models T\) be any countably infinite model.
There is an elementary embedding \(\mathcal A\to\mathcal B\);
it can extend any prescribed correspondence of finite tuples
with the same complete type.
\end{hblemma}

\begin{hbtheorem}{thm:p03-vaught-saturated-uniqueness}{proved=annex:p3.vaught.saturated-uniqueness}{Original Theorem 4.4: uniqueness of the countable saturated model}
Any two countably infinite saturated models of \(T\) are isomorphic.
An elementary correspondence of finite tuples extends to an isomorphism.
\end{hbtheorem}

\begin{hblemma}{thm:p03-vaught-saturated-homogeneity}{proved=annex:p3.vaught.saturated-homogeneity}{Original Lemma 4.5: homogeneous countable saturated models}
Every countably infinite saturated model is homogeneous
in the stated finite complete-type sense.
\end{hblemma}

\begin{hbtheorem}{thm:p03-vaught-saturation-characterization}{proved=annex:p3.vaught.saturation-characterization}{All directions of original Theorem 4.6}
For a countably infinite \(\mathcal A\models T\), the following are
equivalent:
\begin{enumerate}
\item \(\mathcal A\) is saturated.
\item It is \(\aleph_0\)-universal and homogeneous.
\item It is weakly saturated and homogeneous.
\end{enumerate}
The third clause's homogeneity condition is part of the equivalence;
weak realization over the empty set has not been substituted for
the specified finite-parameter extension property.
\end{hbtheorem}

\begin{hbdefinition}{def:p03-vaught-smallness}{The type-counting condition and its stronger neighbor}
The standing theory is \emph{small} when every \(\mathcal P_n(T)\),
\(n<\omega\), is countable. In contrast, \(\omega\)-stability asks
for only countably many complete one-types over every countable
parameter set in a model. The parameter sets and the conditions
are different; smallness does not assert \(\omega\)-stability.
\end{hbdefinition}

\begin{hblemma}{thm:p03-vaught-stability-boundary}{proved=annex:p3.vaught.stability-boundary}{The stronger parameter condition really implies smallness}
For a complete theory in a countable language, \(\omega\)-stability
implies smallness. The converse fails in DLO, as proved below.
This is the parameter-counting comparison for the modern terminology,
not an added hypothesis on original Theorem~4.7.
\end{hblemma}

\begin{hblemma}{thm:p03-vaught-countable-types-atomistic}{proved=annex:p3.vaught.countable-types-atomistic}{The full Boolean-algebra step used in original Corollary 4.8}
If \(\mathcal P_n(T)\) is countable, every nonzero element of
\(\mathbb B_n(T)\) contains an atom.
Indeed an atomless nonzero element would support a binary splitting
tree and produce \(2^{\aleph_0}\) different complete types.
\end{hblemma}

\begin{hbtheorem}{thm:p03-vaught-smallness-equivalences}{proved=annex:p3.vaught.smallness-equivalences}{All four conditions of original Theorem 4.7}
The following are equivalent:
\begin{enumerate}
\item \(T\) is small.
\item It has a countably infinite saturated model.
\item It has an \(\aleph_0\)-universal model.
\item It has a countably infinite weakly saturated model.
\end{enumerate}
The saturated model, when it exists, is unique up to isomorphism
and homogeneous. Saturation retains the finite-parameter meaning
proved above.
\end{hbtheorem}

\begin{hbcorollary}{thm:p03-vaught-universal-prime}{proved=annex:p3.vaught.universal-prime}{Original Corollary 4.8}
If \(T\) has an \(\aleph_0\)-universal model, then it has a prime model.
Equivalently, smallness implies the prime-existence condition here.
The universal model and the prime model are not required to coincide.
\end{hbcorollary}

\subsection{Worked models and the boundaries of the hypotheses}
\label{sec:p03-vaught-examples}

\begin{hbtheorem}{thm:p03-vaught-dlo-example}{proved=annex:p3.vaught.dlo-example}{DLO, finite patterns, named cuts and the stability distinction}
Let \(T_{\rm dlo}\) be dense linear order without endpoints in the
language consisting of \(<\) and identity.
Each complete \(n\)-type over the empty set is isolated by its tuple's
equality and strict-order pattern. Their number is
\[
              \sum_{k=0}^n k!\,\left\{\begin{matrix}n\\k\end{matrix}\right\},
\]
where the braces count partitions of an \(n\)-element set into
\(k\) nonempty blocks; at \(n=0\) the number is one.
The ordered rational model is prime, atomic, saturated and homogeneous.
It is the unique countably infinite model.

Nevertheless there are \(2^{\aleph_0}\) complete one-types over the
countable rational parameter set: distinct real cuts give distinct types.
Thus \(T_{\rm dlo}\) is small but not \(\omega\)-stable.
In the countable language naming every rational, its complete expanded
theory \(T_{\mathbb Q}\) still has the prime atomic model \(\mathbb Q\),
but has \(2^{\aleph_0}\) complete one-types over the empty set of
\emph{that expanded language}. Each irrational cut type is nonprincipal
and omitted by the rational model. Therefore prime existence does not
imply smallness, and the finite-parameter smallness bridge cannot be
extended to arbitrary countable naming.
\end{hbtheorem}

\begin{hbtheorem}{thm:p03-vaught-discrete-example}{proved=annex:p3.vaught.discrete-example}{A source-native prime model and a nonprincipal extension obstruction}
Let \(T_{\rm disc}=\operatorname{Th}(\omega,<)\).
For each \(k<\omega\), the formula saying that \(x\) has exactly
\(k\) predecessors isolates the type of \(k\).
The countably infinite ordered model \((\omega,<)\) is atomic and prime.
The partial type saying that \(x\) has at least \(k\) predecessors
for every \(k\) is consistent and nonprincipal.
Each complete extension of it is nonprincipal and omitted by
that prime model. Thus a prime model need not be weakly saturated.
This works directly with the completed discrete-order example;
its unique finite-position formulas supply the prime-model calculation
without a decision procedure for the full classification of types.
\end{hbtheorem}

\begin{hbtheorem}{thm:p03-vaught-unary-obstruction}{proved=annex:p3.vaught.unary-obstruction}{A countable language with atomless one-variable formulas}
Let \(T_{\rm bits}\) say that every finite Boolean pattern of the
unary predicates \(U_0,U_1,\ldots\) has infinitely many realizations.
Its language has only these predicates and logical identity.
It is a complete satisfiable theory.
Every consistent one-variable formula is atomless.
There are \(2^{\aleph_0}\) complete one-types, one for each
infinite predicate bit sequence, and no prime or atomic model.
The simultaneous construction's atomless branch is indispensable
in this example.
\end{hbtheorem}

\begin{hbtheorem}{thm:p03-vaught-finite-boundary}{proved=annex:p3.vaught.finite-boundary}{Finite complete theories and the cardinal restriction on atomicity}
If a complete theory \(S\) in a countable relational/constant language
has a finite nonempty model \(\mathcal A\), every \(S\)-model is
isomorphic to \(\mathcal A\). Its model is prime, atomic,
homogeneous and saturated over every parameter set within it.
This is a separate finite case; Vaught's denumerable equivalence
was stated with infinite models.

For the complete theory of infinite pure equality, every model
of every infinite cardinality is atomic, whereas its prime models
are precisely the countably infinite ones.
Completeness also matters: in the pure equality language, for the
incomplete theory saying only
that the universe has one or two elements, the one-element model
is atomic in the corresponding formula-algebra sense but cannot
elementarily embed in a two-element model.
\end{hbtheorem}

The source also illustrates its criteria with discrete orders and
real closed fields \cite[p.~315]{Vaught1961}.
For discrete orders it invokes an outside decision procedure for
the full type count and reports the universal homogeneous order
\(\omega+(\omega^*+\omega)\cdot\eta\), where \(\eta\) is the
countable dense order type without endpoints.
For real closed fields it cites the real algebraic field as prime
while the one-type space has continuum many members.
Those additional application claims use the corresponding
order-decision and real-closed-field results. The direct examples
above distinguish prime existence, isolation and smallness.

\subsection{Functional signatures and the broader omission interface}
\label{sec:p03-vaught-language-bridge}

\begin{hblemma}{thm:p03-vaught-functional-bridge}{proved=annex:p3.vaught.functional-bridge}{A complete bridge to countable languages with total functions}
Replace each positive-arity function \(f\) by a relation \(G_f\)
for its graph, with the total and single-valued graph axioms;
zero-arity functions remain individual constants.
There is a cardinality-preserving bijection between the function-language
models of a theory and their graph-language presentations.
Formula translation preserves truth, consistency and free-variable arity,
and has an inverse modulo the graph axioms.
It identifies complete and partial types, principality and isolation,
elementary maps, primeness, atomicity, saturation and homogeneity.
The graph language is countable when the original one is.
Consequently the preceding results transport to countable
function languages with these proved identifications.
The original source presentation remains relational/constant.
\end{hblemma}

\begin{hbtheorem}{thm:p03-vaught-partial-omission}{proved=annex:p3.vaught.partial-omission}{The general countable partial-type omission construction}
Let \(S\) be any satisfiable theory in a countable first-order language,
not necessarily complete, and let \((p_j)_{j\in J}\) be a finite or
countable family of nonprincipal partial types, consistent with \(S\),
in finite tuples of variables.
Then \(S\) has an at most countable model omitting all of them.
It is infinite when \(S\) forces infinitude.
Here principality has the exact consistent-formula meaning of
\cref{def:principal-type}, not merely nonisolation of one selected completion.
Zero-arity types are included.
This generalization uses the witness and finite-refinement construction;
it is distinct from the source's complete-theory simultaneous atom/atomless
theorem and supplies the broader interface of \cref{thm:omitting-types}.
\end{hbtheorem}

\section{Finite Types and Countable Categoricity}
\label{sec:p03-n-a-001-01}

The finite-type criterion connects uniqueness of a countable model to
the number of distinctions expressible about each finite tuple.
The bibliographic title of Ryll-Nardzewski's original article uses
power \(\leq\aleph_0\), including the finite boundary
\cite{RyllNardzewski1959}. Here that boundary is separated from
the infinite-model characterization. The proof develops the
type-realization and extension interfaces of Vaught's formula-algebra
construction (\cref{thm:p03-vaught-finite-formulas-categoricity}),
without identifying this derivation with a different original
author's presentation.

\subsection{The criterion and its hypotheses}
\label{sec:p03-rn-criterion}

\begin{hbconvention}{def:p03-rn-standing}{Complete theories and two meanings of countable categoricity}
In this chapter \(L\) is an at most countable, finite-arity first-order
language with logical identity; constants and function symbols are
allowed. Unless a result explicitly changes the setting, \(T\) is
complete and satisfiable and has an infinite model.
It therefore has only infinite models.
The theory is \(\omega\)-categorical when it has a countably infinite
model and any two countably infinite models are isomorphic.

For the separate finite case, \emph{at-most-countable categoricity}
means existence of a model of size at most \(\aleph_0\) and uniqueness
up to isomorphism among all such models, including finite ones.
These notions agree under the standing infinite-model assumption,
but not without it.

Write \(\bar x=(x_0,\ldots,x_{n-1})\) and let \(S_n(T)\) be the
consistent complete \(n\)-types over the empty set.
Coordinates may repeat and \(n=0\) is included. For
\(\mathcal M\models T\), put
\[
 \operatorname{tp}^{\mathcal M}(\bar a)
       =\{\varphi(\bar x):\mathcal M\models\varphi(\bar a)\}.
\]
The Boolean algebra \(\mathbb B_n(T)\) is the quotient of formulas
in these slots modulo equivalence in every \(T\)-model, as in
\cref{def:p03-vaught-formula-algebra}.
An isolator is consistent with \(T\) and implies every formula
of its complete type; the consistency condition is essential.
The constructions use the compactness and Choice standing of
\cref{thm:fol-compactness}.
\end{hbconvention}

\begin{hblemma}{thm:p03-rn-type-realization}{proved=annex:p3.rn.type-realization}{Every complete finite-tuple type has a countable realization}
For each finite \(n\) and \(p\in S_n(T)\), a countably infinite
\(T\)-model realizes \(p\).
Named coordinates may coincide; no distinctness is imposed unless
the type requires it. The unique zero-type is realized by the empty
tuple. The result includes countable languages with functions.
\end{hblemma}

\begin{hblemma}{thm:p03-rn-isolated-extension}{proved=annex:p3.rn.isolated-extension}{Transfer of an isolated joint type over a fixed tuple}
Suppose \(p\in S_{n+1}(T)\) is isolated by \(\theta(\bar x,y)\),
and \(\bar a\) in a \(T\)-model realizes its first-\(n\)-slot
restriction. Some \(b\) in that same model realizes \(p\)
together with \(\bar a\).
If all empty-parameter finite-tuple types are isolated, every
model realizes all complete finite-tuple types over each finite
parameter set in that model.
\end{hblemma}

\begin{hbtheorem}{thm:p03-rn-characterization}{proved=annex:p3.rn.characterization}{The complete finite-type characterization}
Under \cref{def:p03-rn-standing}, the following are equivalent:
\begin{enumerate}
\item \(T\) is \(\omega\)-categorical.
\item \(S_n(T)\) is finite for every \(n<\omega\).
\item Every \(p\in S_n(T)\), for every finite \(n\), is isolated.
\item \(\mathbb B_n(T)\) is finite for every finite \(n\).
\item Every model of \(T\) is atomic.
\item Every model of \(T\) is \(\aleph_0\)-saturated:
      it realizes types over finite parameter sets.
\item Every countably infinite model of \(T\) is
      \(\aleph_0\)-saturated.
\item \(T\) has a countably infinite model that is both prime
      and \(\aleph_0\)-saturated.
\end{enumerate}
The finite number \(|S_n(T)|\) may depend on \(n\);
no single bound for all tuple lengths is required.
There is no finite-signature or function-free restriction
on this general theorem.
\end{hbtheorem}

The uniqueness-to-isolation obstruction is a type realized
in one countable model and omitted in another. In the reverse
direction, isolation permits an extension over each already
matched finite tuple. Alternating extensions in two enumerated
models gives an onto elementary map. Neither mere density
of isolated types nor a single successful tuple extension
would suffice for this criterion.

\begin{hbcorollary}{thm:p03-rn-orbits}{proved=annex:p3.rn.orbits}{Types are the finite-tuple automorphism orbits}
If \(T\) satisfies the equivalent conditions and
\(\mathcal M\models T\) is countably infinite, two finite
tuples of the same length are in the same
\(\operatorname{Aut}(\mathcal M)\)-orbit exactly when
they have the same complete type.
Every type in \(S_n(T)\) is realized in \(\mathcal M\),
so there are exactly \(|S_n(T)|\) orbits on its \(n\)-tuples.
This includes the one orbit on zero-tuples.
The construction extends the prescribed tuple map.
\end{hbcorollary}

\subsection{Finite, incomplete and uncountable-language boundaries}
\label{sec:p03-rn-boundaries}

\begin{hbtheorem}{thm:p03-rn-finite-boundary}{proved=annex:p3.rn.finite-boundary}{The finite complete-theory alternative}
Let \(T\) be complete and satisfiable in an at most countable
language and have a nonempty finite model of size \(m\).
Every \(T\)-model is isomorphic to it.
Thus \(T\) is at-most-countable categorical, but not
\(\omega\)-categorical under the existence convention above.
For every finite \(n\),
\[
 |S_n(T)|\leq m^n,\qquad
 |\mathbb B_n(T)|\leq2^{m^n}.
\]
Every complete type is isolated and realized in this model,
which is atomic, prime and saturated over every parameter
set in its finite domain. A finite-type criterion alone
does not imply existence of an infinite model.
\end{hbtheorem}

\begin{hbtheorem}{thm:p03-rn-incomplete-boundary}{proved=annex:p3.rn.incomplete-boundary}{Infinite-model uniqueness does not replace completeness}
In the language with only logical identity, the empty theory
has isomorphic countably infinite models but infinitely many
complete one-types: each positive finite exact-cardinality
sentence belongs to a different type.
It also has models of every positive finite size.
Thus uniqueness at the countably infinite cardinal, without
completeness or an assumption that all models are infinite,
does not imply the finite-type criterion.
\end{hbtheorem}

\begin{hbtheorem}{thm:p03-rn-uncountable-language}{proved=annex:p3.rn.uncountable-language}{A genuine countable-language boundary}
There is a complete theory in an uncountable language with
exactly one at most countable model, which is infinite,
but infinitely many complete one-types.
Take the complete theory of \(\omega\) with constants \(c_i\)
for all \(i<\omega\) and a unary function \(f_g\) for every
\(g:\omega\to\omega\), interpreted by \(g\).
Every model containing an element outside all the named
constants has at least \(2^{\aleph_0}\) elements.
The language's countability is an actual hypothesis in
the omission-based converse.
\end{hbtheorem}

\begin{hbconvention}{def:p03-rn-parameter-boundary}{Finite parameters are not countable parameters}
The criterion's \(\aleph_0\)-saturation concerns finite
parameter sets, not arbitrary countable ones.
In any countably infinite model, the consistent type asking
for an element different from every member of its domain
is omitted over that whole domain.
The type count is over the empty set of a fixed language;
infinitely many new names can destroy it.
The rationally named dense order and its irrational-cut
types in \cref{thm:p03-vaught-dlo-example} give the explicit
example.
\end{hbconvention}

\subsection{Dense linear order: the full type calculation}
\label{sec:p03-rn-dlo}

\begin{hbdefinition}{def:p03-rn-dlo-data}{Dense order without endpoints and finite order patterns}
The language here has just \(<\) and logical identity:
no named points, extra predicates or function symbols.
Let \(T_{\rm dlo}\) consist of the strict linear-order axioms,
density
\[
 \forall x\forall y\bigl(x<y\to
                     \exists z(x<z\land z<y)\bigr),
\]
and the absence of both endpoints
\[
 \forall x\exists y(y<x),\qquad
 \forall x\exists y(x<y).
\]
An order pattern is an ordered partition \((B_1,\ldots,B_k)\)
of \(\{0,\ldots,n-1\}\) into nonempty blocks.
Its formula \(\delta_{(B_1,\ldots,B_k)}(\bar x)\)
requires equal coordinates within each block and
\(x_i<x_j\) whenever \(i\in B_r\), \(j\in B_s\), \(r<s\).
There is one empty pattern for \(n=0\), with formula \(\top\).
The number \({n\brace k}\) counts partitions into exactly
\(k\) nonempty unordered blocks; \({0\brace0}=1\).
\end{hbdefinition}

\begin{hblemma}{thm:p03-rn-dlo-extension}{proved=annex:p3.rn.dlo-extension}{Every finite order map extends in both directions}
An order-preserving bijection between finite subsets of two
\(T_{\rm dlo}\)-models extends to any further specified
element of either domain.
An old element uses its existing mate.
A new one is matched below the first old point, above the last,
or between the appropriate adjacent old points.
With an empty initial subset any target point suffices.
Each open interval and end ray is infinite, so a fresh target
can be selected when needed.
No countability hypothesis is used in this finite extension.
\end{hblemma}

\begin{hbtheorem}{thm:p03-rn-dlo-pattern-truth}{proved=annex:p3.rn.dlo-pattern-truth}{Order patterns decide every formula}
Tuples in any two \(T_{\rm dlo}\)-models with the same order
pattern satisfy the same first-order formulas in their slots.
Thus \(T_{\rm dlo}\) is complete, and each formula is
equivalent modulo it to the finite disjunction of the order
patterns making it true.
This proves quantifier elimination without assuming it.
Every pattern is realized in every model.
\end{hbtheorem}

\begin{hbcorollary}{thm:p03-rn-dlo-type-count}{proved=annex:p3.rn.dlo-type-count}{Ordered partitions are exactly the complete types}
The pattern formulas isolate exactly the complete types.
Therefore
\[
 |S_n(T_{\rm dlo})|=\sum_{k=0}^n k!\,{n\brace k},
 \qquad
 |\mathbb B_n(T_{\rm dlo})|=2^{\,|S_n(T_{\rm dlo})|}.
\]
The type counts for \(n=0,1,2,3\) are \(1,1,3,13\).
The two-coordinate types are equality and the two strict
orderings; the thirteen three-coordinate types are one
all-equal pattern, six two-block patterns and six three-block
patterns.
\end{hbcorollary}

\begin{hbtheorem}{thm:p03-rn-dlo-uniqueness}{proved=annex:p3.rn.dlo-uniqueness}{The rational model and prescribed-map back-and-forth}
The ordered rationals form a countably infinite
\(T_{\rm dlo}\)-model. Every countably infinite dense linear
order without endpoints is isomorphic to it by alternating
finite order extensions.
The construction extends any prescribed finite order map
between two such models.
Consequently the rational model is homogeneous, its tuple
orbits are precisely the ordered partitions, and
\(T_{\rm dlo}\) is \(\omega\)-categorical.
Density without both endpoint clauses would not justify
the end-ray extension steps.
\end{hbtheorem}

\subsection{The corrected finite-relational application}
\label{sec:p03-rn-finite-relational}

The general theorem allows functions and infinite countable
signatures. A useful \emph{application} has stronger hypotheses:
Hodges's author correction to the finite-signature application
at p.~352 adds the absence of function symbols
\cite[p.~3, correction to p.~352]{Hodges1993Corrigenda2011}.
The following argument makes that restriction's role explicit.

\begin{hbtheorem}{thm:p03-rn-relational-application}{proved=annex:p3.rn.relational-application}{Finite relational homogeneity implies countable categoricity}
Let \(\mathcal M\) be countably infinite in a finite,
finite-arity, purely relational language without functions
or constants. Suppose every isomorphism between finite
induced substructures extends to an automorphism.
Then \(\operatorname{Th}(\mathcal M)\) has quantifier
elimination and is \(\omega\)-categorical.
For each fixed \(n\), the finitely many atomic relation
and equality patterns give a finite list of complete types.
Finiteness of this list and extendibility of the finite maps
are separate steps.
\end{hbtheorem}

\begin{hbtheorem}{thm:p03-rn-function-boundary}{proved=annex:p3.rn.function-boundary}{Finite function signature does not imply a finite pattern list}
Let \(\mathcal Z=(\mathbb Z,s,t)\), with
\(s(a)=a+1\) and \(t(a)=a-1\).
Every nonempty finitely generated substructure is the whole
structure; an isomorphism between such substructures is
already an automorphism.
Nevertheless its complete theory is not
\(\omega\)-categorical: the pairs \((0,k)\), \(k\geq0\),
have distinct two-types, distinguished by \(y=s^k(x)\).
There is no nonempty finite closed substructure.
This example does not satisfy the finite induced-substructure
hypothesis; it shows why finite function signatures cannot
silently be treated as finite relational pattern calculations.
\end{hbtheorem}

The source identities have distinct extents. The selected
Engeler 1959 entry is a Notices announcement at p.~161,
not a longer proof article \cite{Engeler1959}.
The selected Svenonius categoricity article is Theoria 25,
pp.~82--94 \cite{Svenonius1959}, distinct from his same-year
permutation article at pp.~173--178.
The Hodges correction concerns the bounded application,
not a new function-free hypothesis on the general criterion.


\section{Countable saturated models and Vaught's problem}
\label{sec:p03-countable-models}

Call a complete theory in a countable language \emph{small} when every
space \(S_n(T)\), \(n<\omega\), is countable.  This adjective is needed here
to distinguish the exact type-counting condition for a countable saturated
model from the stronger condition of \(\omega\)-stability introduced below.
It is a property of a complete theory, not of one presentation, and it is used
only in the following existence theorem.  In particular, DLO is small but is
not \(\omega\)-stable.

\begin{hbtheorem}{thm:p03-countable-saturated}{proved=annex:p3.vaught.overview-saturated}{Countable saturated-model package}
Let \(T\) be a complete theory in a countable language with an infinite
model.
\begin{enumerate}
  \item The theory \(T\) has a countable saturated model if and only if it is
        small.
  \item Any two countable saturated models of \(T\) are isomorphic.
\end{enumerate}
Here ``saturated'' uses the handbook convention: a countable model realizes
every complete finite-tuple type over every finite parameter set.
\end{hbtheorem}

The full nested elementary-chain construction is
\cref{thm:p03-vaught-saturated-existence};
\cref{thm:p03-vaught-smallness-equivalences} supplies both directions of
clause~1, and \cref{thm:p03-vaught-saturated-uniqueness} supplies clause~2.
The parameter correspondence and function-language bridge retain precisely
the finite-parameter saturation used here. The independent presentation in
\cite{Marker2002} remains separately attributed.

For a complete theory \(T\) in a countable language, write
\(I(T,\aleph_0)\) for the number of isomorphism types of countable models of
\(T\).  This notation is needed to state the countable-model problem, stands
for a cardinal-valued metalinguistic invariant of \(T\), and is not the domain
of a model or a new categoricity relation.  It differs from
\(\aleph_0\)-categoricity, which is precisely the equation
\(I(T,\aleph_0)=1\), and it is used in the next theorem and conjecture.

\begin{hbtheorem}{thm:p03-vaught-never-two}{proved=annex:p03.vaught-two.overview}{Vaught's never-two theorem}
No complete theory in a countable language has exactly two nonisomorphic
countable models: \(I(T,\aleph_0)\ne2\).
\end{hbtheorem}

Vaught's conjecture asks whether every complete theory in a countable
language satisfies
\[
 I(T,\aleph_0)\leq\aleph_0
 \qquad\text{or}\qquad
 I(T,\aleph_0)=2^{\aleph_0}.
\]
It remains open in this generality.  The never-two theorem is therefore not a
solution of the conjecture, and Morley's uncountable categoricity theorem does
not decide arbitrary countable-model spectra.  The Baldwin--Lachlan theorem
below records an important class for which the conjectured dichotomy is known
\cite{Vaught1961,BaldwinLachlan1971,Marker2024}.

\section{Vaught's never-two theorem and countable-model problems}
\label{sec:p03-vaught-never-two-presentation}

Vaught's original Section~6 answers a question of Raphael Robinson. The
introduction of the paper announces that Section~6 shows that a complete
theory cannot have exactly two nonisomorphic denumerable models, and that
this answers Robinson's question \cite[p.~304]{Vaught1961}. Section~6 then
states the question. Robinson had observed that the complete theory of an
equivalence relation with exactly two classes, both infinite, has exactly two
nonisomorphic models of power \(\aleph_1\), and he asked whether some
complete theory has exactly two nonisomorphic models of power \(\aleph_0\).
A footnote dates the question to a conversation in 1957 among several
people, Vaught included \cite[p.~318 and n.~23]{Vaught1961}. Vaught's
answer combines examples with a restriction: Theorem~6.1 says that a
complete theory with exactly \(n\) nonisomorphic denumerable models exists
if and only if \(n\ne2\) \cite[pp.~318--320]{Vaught1961}. The paper's letters
\(i,\ldots,r\) range over the natural numbers \(0,1,\ldots\)
\cite[p.~304, n.~5]{Vaught1961}, so the theorem covers every \(n\in\omega\),
including \(n=0\); Theorem~\ref{thm:p03-w03-p03-016-theorem-6-1} states
it in that form. In Vaught's usage a complete theory is consistent
\cite[p.~303, n.~1]{Vaught1961}. Throughout this presentation a theory is
consistent, and countable includes finite unless ``denumerable'' is specified.
All languages here are countable. Until the final graph translation in the
never-two proof, the language has only relation symbols and constants, as in
the original presentation.
The graph translation described in the proof transfers the negative result
to languages with function symbols.

\subsection{Type counting and the named-type theory}
\label{sec:p03-vaught-never-two-types}

\begin{hbdefinition}{def:p03-vaught-spectrum-notation}{Spectra, formula algebras and named types}
For a complete first-order theory \(T\), \(I(T,\aleph_0)\) counts its
countable models up to isomorphism. Write \(S_m(T)\) for its complete
\(m\)-types and \(F_m(T)\) for formulas in the first \(m\) free variables
modulo equivalence in \(T\). Boolean conjunction, disjunction and negation
make \(F_m(T)\) a Boolean algebra. Vaught's \(\mathfrak P_m(T)\), the set
of dual prime ideals of this algebra, corresponds to \(S_m(T)\): a type
contains exactly the formulas in its associated dual prime ideal.

For \(p\in S_m(T)\) and fresh constants \(\bar c=(c_0,\ldots,c_{m-1})\),
put \(T_p=T\cup\{\varphi(\bar c):\varphi\in p\}\).
This is a complete theory in the expanded language. Its models are precisely
pointed models \((M,\bar a)\) with \(\operatorname{tp}_M(\bar a)=p\).
Inclusion of the original language induces a map from \(F_k(T)\) to
\(F_k(T_p)\); injectivity is a theorem below, not part of this definition.
\end{hbdefinition}

The complete-type and isolation conventions are those of
Definition~\ref{def:principal-type}. Saturation here concerns finite tuples
over finite parameter sets, as in Theorem~\ref{thm:p03-countable-saturated}.
The original existence and back-and-forth constructions are supplied in
Theorems~\ref{thm:p03-vaught-saturated-existence} and
\ref{thm:p03-vaught-saturated-embedding}.

\begin{hblemma}{thm:p03-vaught-few-models-types}{proved=annex:p03.vaught-two.types}{Few countable models give small type spaces}
If \(T\) has an infinite model and \(I(T,\aleph_0)\leq\aleph_0\), every
\(S_m(T)\) is countable. In every countable type space, isolated types are
dense. A type space in which every type is isolated is finite; conversely,
every type in a finite type space is isolated.
\end{hblemma}

Thus the hypothesis of exactly two models supplies smallness. It does not
supply \(\omega\)-stability, which counts types over all countable parameter
sets. The original prime/atomic construction
(Theorem~\ref{thm:p03-vaught-prime-atomic}) and saturation construction
operate under their displayed hypotheses; neither adjective may be exchanged
for the other in this argument.

\begin{hblemma}{thm:p03-vaught-named-finiteness}{proved=annex:p03.vaught-two.named}{Categoricity after naming a complete type}
Let \(T\) be complete with infinite models. If \(T_p\) has exactly one
denumerable model up to isomorphism, every \(F_k(T_p)\) is finite.
Moreover, \(F_k(T)\to F_k(T_p)\) is injective for every \(k\).
Consequently, if some \(S_m(T)\) contains a nonisolated type, \(T_p\)
cannot be denumerably categorical for any \(p\).
\end{hblemma}

\begin{hbtheorem}{thm:p03-vaught-never-two-original}{proved=annex:p03.vaught-two.never}{Vaught's never-two theorem}
A complete first-order theory in a countable language cannot have exactly
two nonisomorphic countable models. In particular,
\(I(T,\aleph_0)\ne2\).
\end{hbtheorem}

The proof names a realization of a nonisolated type. A model omitting that
type is one of the hypothetical two isomorphism types. Every model realizing
it must therefore be the other, saturated type. Back-and-forth then makes
\(T_p\) categorical, contradicting
Lemma~\ref{thm:p03-vaught-named-finiteness}. Finite models and the absence
of a nonisolated type are separate cases in the annex.

\subsection{The examples surrounding the theorem}
\label{sec:p03-vaught-finite-count-examples}

\begin{hbdefinition}{def:p03-vaught-colored-orders}{Ehrenfeucht's orders and their finite-language coding}
Fix a positive integer \(r\). Let \(T_{r+2}\) assert that \(<\) is a dense
linear order without endpoints, the unary predicates
\(U_0,\ldots,U_{r-1}\) partition its domain and are each dense, and
\(c_0<c_1<\cdots\), with every \(c_i\in U_0\).
For \(r=1\), discard the redundant predicate and obtain the original
three-model theory \(T_3\).

Its finite-language coding \(T'_{r+2}\) has an equivalence relation \(R\),
an order \(<\) constant on \(R\)-classes, and class-constant predicates
\(U_j\). On the quotient the order is dense without endpoints and each
color is dense. For every positive integer \(i\), exactly one class
\(C_i\) has size \(i\); these classes have color \(U_0\) and satisfy
\(C_i<C_{i+1}\). All other classes are infinite. ``Class has size \(i\)''
is its ordinary first-order formula with \(i\) distinct members and a
universal assertion that these exhaust the class. The notation \(C_i\)
denotes a definable quotient point, not an added constant symbol.
For \(r=1\), this is Vaught's displayed language \((<,R)\).
\end{hbdefinition}

\begin{hbtheorem}{thm:p03-vaught-finite-counts}{proved=annex:p03.vaught-two.examples}{All positive finite counts except two}
For every positive integer \(n\), a complete first-order theory has exactly
\(n\) nonisomorphic denumerable models if and only if \(n\ne2\).
For \(n=1\), dense linear order without endpoints suffices. For \(n=r+2\),
both \(T_{r+2}\) and \(T'_{r+2}\) are complete and have exactly \(r+2\)
denumerable models. The latter has a finite language.
\end{hbtheorem}

For the named-constant theory the possibilities are: the constants are
cofinal; they have upper bounds but no least upper bound in the domain; or
their least upper bound exists and has one of the \(r\) colors. The quotient
points \(C_1,C_2,\ldots\) give the same alternatives for the coding.
The upper-bound condition is the infinite conjunction
\(\bigwedge_i c_i<x\); it is not an additional first-order predicate.
The annex proves completeness as well as the isomorphism classification.

Vaught's own justifications in this part of the proof are short citations.
That the theory \(T_1\) of dense linear order without endpoints is complete
with exactly one denumerable model up to isomorphism is credited to Cantor's
theorem \cite[p.~318]{Vaught1961}. For \(T_n\) with \(n\ge4\), which is
\(T_{r+2}\) above with \(r=n-2\), he invokes a theorem of
Skolem (his reference~[14]): any two denumerable models of \(T_n\) are
isomorphic once their lists of distinguished elements are ignored. He states
that completeness and the list of isomorphism types are easily seen from this
fact \cite[p.~319]{Vaught1961}. The colored back-and-forth argument at the
start of the annex proof of Theorem~\ref{thm:p03-vaught-finite-counts} is a
proof of Skolem's fact as Vaught uses it, and its uncolored case is Cantor's
theorem. Clause~1 of Theorem~\ref{thm:p03-w03-p03-016-source-routes}
derives the completeness of \(T_{r+2}\) from Skolem's fact in the way the
source indicates.

For the finite-language theory Vaught sketches a different completeness
argument, granted that the theory has only three denumerable models
\cite[pp.~319--320]{Vaught1961}. Start from a denumerable model of \(T'_3\)
in which the finite classes are cofinal. Adding a constant above all of
them and applying the completeness theorem yields an elementarily equivalent
denumerable model of one of the other two kinds. Further applications of the
completeness theorem, together with the union of an elementarily increasing
family, show that a model in which the finite classes have a limit is
elementarily equivalent to one in which they neither have a limit nor are
cofinal. Completeness then follows from the L\"owenheim--Skolem theorem.
Clause~2 of Theorem~\ref{thm:p03-w03-p03-016-source-routes} carries out this
route in full for every \(T'_{r+2}\); the quantifier-rank game in the annex
proof of Theorem~\ref{thm:p03-vaught-finite-counts} is an independent second
proof.

\begin{hbtheorem}{thm:p03-w03-p03-016-source-routes}{proved=annex:p03-w03-p03-016-source-routes}{Completeness by Vaught's two routes}
Let \(r\) be a positive integer.
\begin{enumerate}
  \item (Skolem's fact.) Every sentence of the language of \(T_{r+2}\)
        mentions only \(c_0,\ldots,c_k\) for some \(k\). The reducts of any
        two denumerable models of \(T_{r+2}\) to
        \(<,U_0,\ldots,U_{r-1},c_0,\ldots,c_k\) are isomorphic. Hence
        \(T_{r+2}\) is complete.
  \item (The finite-language sketch.) Let \(A\) be a denumerable model of
        \(T'_{r+2}\) in which the finite classes are cofinal. Some
        denumerable model of \(T'_{r+2}\) in which the finite classes are
        bounded is elementarily equivalent to \(A\). Every denumerable model
        of \(T'_{r+2}\) in which the finite classes have a least upper bound
        is elementarily equivalent to a denumerable model in which they are
        bounded and have no least upper bound. Consequently all \(r+2\)
        denumerable isomorphism types are elementarily equivalent, and
        \(T'_{r+2}\) is complete.
\end{enumerate}
\end{hbtheorem}

\begin{hbtheorem}{thm:p03-w03-p03-016-theorem-6-1}{proved=annex:p03-w03-p03-016-theorem-6-1}{Theorem~6.1 over all natural numbers}
For every \(n\in\omega\) there is a complete theory, in a countable language
whose nonlogical symbols are relation symbols and individual constants, with
exactly \(n\) nonisomorphic denumerable models, if and only if \(n\ne2\).
For \(n=0\) the complete theory of a one-element structure is a witness.
Every complete theory in a countable language with an infinite model has at
least one denumerable model.
\end{hbtheorem}

The case \(n=0\) is not discussed in Vaught's text. Its witnesses have only
finite models, so they lie outside the standing assumption, adopted in
\S\,1, that \(T\) is a complete theory with infinite models
\cite[p.~306]{Vaught1961}. The last sentence of the theorem shows that,
inside that assumption, the finite counts which occur are exactly the
positive \(n\ne2\).

\begin{hbtheorem}{thm:p03-vaught-equivalence-example}{proved=annex:p03.vaught-two.equivalence}{Two models at an uncountable cardinal}
The theory asserting exactly two infinite equivalence classes is complete.
It has one denumerable model, but exactly two models of size \(\aleph_1\):
the unordered class-size pairs are
\(\{\aleph_0,\aleph_1\}\) and \(\{\aleph_1,\aleph_1\}\).
\end{hbtheorem}

This is the theory from which Robinson's question starts
\cite[p.~318]{Vaught1961}. Its two models occur at the cardinal
\(\aleph_1\), not at the countable cardinal. The never-two theorem answers
the question negatively: no complete theory has exactly two nonisomorphic
denumerable models. If finite models are counted as countable, a complete
theory with a finite model has exactly one countable isomorphism type, and a
complete theory with infinite models has a denumerable one, so that count is
never \(0\). Counting denumerable models, as Vaught does, the value \(0\)
occurs, by Theorem~\ref{thm:p03-w03-p03-016-theorem-6-1}.

\subsection{The original questions and the later formulations}
\label{sec:p03-vaught-original-problem}

Immediately after the proof of Theorem~6.1, Vaught records a conjecture
which he calls unresolved and which, he says, appears to have been made by
a number of people. In his words: ``A theory which is
\(\aleph_{1+\alpha}\)-categorical but not \(\aleph_0\)-categorical has
exactly \(\aleph_0\) non-isomorphic denumerable models.''
\cite[p.~320]{Vaught1961}. The count is of denumerable models, so finite
models play no part in it, and the hypothesis is categoricity in some
uncountable power \(\aleph_{1+\alpha}\). Theories here have at most
\(\aleph_0\) nonlogical constants, each a relation symbol or an individual
constant \cite[p.~304]{Vaught1961}. From a point in \S\,1 onward the paper
assumes that the theory \(T\) under discussion is complete with infinite
models, while remarking that one or both assumptions can often be dropped
\cite[p.~306]{Vaught1961}. The conjecture itself says only ``a theory''.
He then poses a second problem: whether it can be proved, without the
continuum hypothesis, that some complete theory has exactly \(\aleph_1\)
nonisomorphic denumerable models \cite[p.~320]{Vaught1961}. These are the
actual original questions. The cardinal dichotomy customarily called
Vaught's conjecture is formulated below; it does not appear in this form in
the 1961 text. The introduction of the selected 2026 status source
attributes the countable-or-continuum statement to this paper
\cite[p.~1]{GonzalezRosseggerTuretsky2026Status}; the 1961 text itself
displays only the two questions just quoted.

The categoricity conjecture is a statement about uncountably categorical
theories, and it is distinct from the never-two theorem: the never-two
theorem excludes one finite count for every complete theory and says nothing
about which theories have \(\aleph_0\) denumerable models. For a complete
theory in a countable language, the later Baldwin--Lachlan refinement
(Theorem~\ref{thm:p03-baldwin-lachlan}, a cited result) gives
\(I(T,\aleph_0)\in\{1,\aleph_0\}\) under categoricity in one uncountable
cardinal. Such a theory has no finite models, so this count is a count of
denumerable models. If \(T\) is not \(\aleph_0\)-categorical, the value
\(1\) is excluded and the count is \(\aleph_0\), which is Vaught's
conjecture for complete theories.

\begin{hbdefinition}{def:p03-vaught-variants}{Cardinal, perfect-set and infinitary claims}
The \emph{first-order cardinal conjecture} says, for every complete theory
in a countable language,
\[
 I(T,\aleph_0)\leq\aleph_0\quad\text{or}\quad
 I(T,\aleph_0)=2^{\aleph_0}.
\]
The \emph{perfect-set formulation} asks instead for a perfect set of
pairwise nonisomorphic models whenever there are uncountably many
isomorphism types. Models with domain \(\omega\) are coded by the product
space of their relation, function and constant interpretations; finite
domains have the analogous coding components. Perfect means nonempty,
closed and without isolated points in the resulting coding space.

The \emph{infinitary cardinal conjecture} makes the same counting assertion
for the countable models of each \(L_{\omega_1\omega}\)-sentence
\(\sigma\). This logic permits countable conjunctions and disjunctions,
with only finitely many free variables in such a connective and finite
quantifier strings. It is a wider axiomatization class than a complete
first-order theory. The perfect-set question can likewise be asked in this
infinitary class.
\end{hbdefinition}

\begin{hblemma}{thm:p03-vaught-cardinal-boundaries}{proved=annex:p03.vaught-two.variants}{What cardinal counting does and does not assert}
In a countable language there are at most \(2^{\aleph_0}\) countable
isomorphism types. A perfect set of pairwise nonisomorphic models supplies
\(2^{\aleph_0}\) types. Under CH, the bare countable-or-continuum cardinal
dichotomy follows from this upper bound for any class of countable
structures. That counting argument supplies no perfect set.
There is an \(L_{\omega_1\omega}\)-sentence with exactly two denumerable
models whose first-order reducts both satisfy the complete theory \(T_3\).
\end{hblemma}

Every countable first-order theory can be conjoined into one infinitary
sentence, so the infinitary cardinal claim entails the first-order one.
The reverse implication is not supplied by this translation. Likewise,
cardinality alone does not supply the topological witnesses in the perfect-set
formulation. Morley's separate 1970 source treatment concerns its precise
infinitary threshold; the never-two proof does not establish that theorem.

\subsection{Scott ranks and the dated status}
\label{sec:p03-vaught-status-2026}

\begin{hbdefinition}{def:p03-vaught-scott-comparison}{Ranks used in the status comparison}
A Scott sentence of a countable structure characterizes it up to isomorphism
among countable structures. In the selected 2026 convention,
\(SR(M)\) is the least \(\alpha\) for which such a sentence is
\(\Pi^{\mathrm{in}}_{\alpha+1}\), and \(SR_p(M)\) is the least
\(\alpha\) for which one is \(\Sigma^{\mathrm{in}}_{\alpha+2}\).
Here \(\Sigma^{\mathrm{in}}_0=\Pi^{\mathrm{in}}_0\) are finitary
quantifier-free formulas; higher \(\Sigma\)-classes are countable
disjunctions of existentially quantified lower \(\Pi\)-formulas, and the
\(\Pi\)-classes are the dual conjunctions of universally quantified lower
\(\Sigma\)-formulas. Parameterized rank permits one fixed finite tuple of
named parameters; it is not interchangeable with unparameterized rank.

Write \(\operatorname{vo}(\sigma)\) for the Vaught ordinal: the least
ordinal \(\beta\), if one exists, such that either continuum many
\(\Pi^{\mathrm{in}}_\beta\)-types are realized among the models of
\(\sigma\), or there are only countably many models of \(\sigma\) and they
all have Scott rank less than \(\beta\); otherwise use \(\omega_1\)
\cite[p.~6]{GonzalezRosseggerTuretsky2026Status}. Here the
\(\Pi^{\mathrm{in}}_\beta\)-type of a tuple is the set of
\(\Pi^{\mathrm{in}}_\beta\)-formulas it satisfies, and, as Scott ranks
require, the models are the countable models of \(\sigma\). Since negation
exchanges \(\Pi^{\mathrm{in}}_\beta\) and \(\Sigma^{\mathrm{in}}_\beta\)
up to logical equivalence, the \(\Pi^{\mathrm{in}}_\beta\)-type and the
\(\Sigma^{\mathrm{in}}_\beta\)-type of a tuple determine each other.

The \emph{\(\omega\)-Vaught conjecture}, in the formulation the selected
source attributes to Gonzalez and Montalb\'an, states that every theory
has a countable Vaught ordinal and, more strongly, that every
\(\Pi^{\mathrm{in}}_\alpha\)-theory \(\sigma\) satisfies
\(\operatorname{vo}(\sigma)\leq\alpha+\omega\)
\cite[p.~4]{GonzalezRosseggerTuretsky2026Status}. In the infinitary setting
the source describes it as the translation of Martin's conjecture, a
model-theoretic strengthening of Vaught's conjecture
\cite[p.~4]{GonzalezRosseggerTuretsky2026Status}.
\end{hbdefinition}

The selected source states that the existence of a countable Vaught ordinal
for \(\sigma\) is equivalent to \(\sigma\) satisfying Vaught's conjecture,
in its infinitary formulation, and that a counterexample has Vaught ordinal
\(\omega_1\) \cite[pp.~2, 6]{GonzalezRosseggerTuretsky2026Status}. The
\(\omega\)-Vaught conjecture therefore strengthens the infinitary conjecture:
it keeps the countability of every Vaught ordinal and adds the bound
\(\alpha+\omega\) determined by the syntactic complexity of the sentence.
Theorem~\ref{thm:p03-w03-p03-016-vo-cardinal} proves, in ZFC, the direction
from a countable Vaught ordinal to the cardinal dichotomy, and explains why
the reverse direction cannot be read as a statement about the bare cardinal
dichotomy alone once the continuum hypothesis is assumed.

\paragraph{Edition status, 9 September 2026.}
The selected evidence is Gonzalez, Rossegger and Turetsky,
arXiv:2606.15205v2, dated 26 August 2026. Its introduction treats the general
conjecture as open and subsequently uses the infinitary formulation.
Corollary~2.11 concerns hypothetical \(\Pi^{\mathrm{in}}_2\)
counterexamples: they must have two models at every positive countable
parameterized Scott rank. Theorem~3.6 gives a different conclusion: a minimal
counterexample has exactly one model at each rank in an unbounded set of
unparameterized Scott ranks. These statements do not contradict one another.
Corollary~4.10 proves the \(\omega\)-Vaught bound for infinitary theories
with finitely many models, not for arbitrary theories. This restricted
result and the never-two theorem leave the general conjecture unresolved
\cite{GonzalezRosseggerTuretsky2026Status}.

The version record reports added assumptions in Section~3. The displayed
status therefore belongs to v2 and to the edition's historical date; it does
not silently combine v1 statements, later versions, or the three different
counting questions.

\section{Countable-Model Spectra and Infinitary Isomorphism}
\label{sec:p03-n-a-004-02}

The number of isomorphism types is not the number of labeled copies.
Counting countable models therefore involves both a coding space
and an equivalence relation on it. Morley's 1970 threshold concerns
more than \(\aleph_1\) isomorphism types, not merely an uncountable
set of codes \cite{Morley1970CountableModels}.
The construction below supplies the infinitary isomorphism and
coding interfaces used in such spectrum arguments.

\subsection{Codes and the threshold}
\label{sec:p03-morley-codes}

\begin{hbdefinition}{def:p03-morley-coding}{The fixed-domain coding space}
Let \(L\) be an at most countable first-order language of
finite-arity relations, functions and constants, with logical
identity. Structures are nonempty. Write \(X_L\) for its
structures with domain \(\omega\), coded as
\[
 X_L=
 \prod_{R\in L}2^{\,\omega^{\operatorname{arity}(R)}}
 \ \times\
 \prod_{f\in L}\omega^{\,\omega^{\operatorname{arity}(f)}}
 \ \times\
 \prod_{c\in L}\omega .
\]
Positive-arity functions are in the second product;
individual constants are in the third. Nullary relations have
one Boolean coordinate. Each discrete coordinate has its
discrete topology and the whole space has the product topology.
Codes record all interpretations; logical identity is actual equality.
Let \(X_{L,m}\) be the analogous space on
\(\{0,\ldots,m-1\}\), for \(m\geq1\).
The topological sum \(X_L^{\leq\omega}\) of these spaces and
\(X_L\) includes all nonempty at most countable structures.

For a class \(K\subseteq X_L\), put
\(I(K)=|K/{\cong}|\). For a sentence \(\sigma\) or theory \(T\),
\(K_\sigma\) or \(K_T\) is its model class on \(\omega\).
These countably infinite-domain counts are kept distinct
from the all-at-most-countable convention of
\cref{def:p03-vaught-spectrum-notation}.
\end{hbdefinition}

\begin{hblemma}{thm:p03-morley-code-space}{proved=annex:p3.morley.codes}{The coding spaces are Polish and have the continuum upper bound}
The spaces \(X_L\), \(X_{L,m}\) and \(X_L^{\leq\omega}\)
are Polish: separable and completely metrizable.
Their cardinalities are at most \(2^{\aleph_0}\).
Consequently any class of at most countable \(L\)-structures
has at most \(2^{\aleph_0}\) isomorphism types.
For an infinite language, even a fixed finite domain can
have continuum many nonisomorphic structures; finite size
alone is not a countable bound on the number of theories.
\end{hblemma}

\begin{hbtheorem}{thm:p03-morley-firstorder-threshold}{cited=Morley1970CountableModels}{Morley's first-order threshold}
For a complete satisfiable first-order theory \(T\) in a
countable language with infinite models,
\[
 I(K_T)>\aleph_1
       \quad\Longrightarrow\quad I(K_T)=2^{\aleph_0}.
\]
This is a first-order instance of the 1970 threshold.
It leaves the case \(I(K_T)=\aleph_1<2^{\aleph_0}\)
untouched; it does not replace \(>\aleph_1\) by
\(>\aleph_0\).
\end{hbtheorem}

Morley's published introduction also announces extensions to
\(L_{\omega_1,\omega}\) and pseudo-axiomatic classes in that
logic \cite[p.~14]{Morley1970CountableModels}.
An arbitrary collection of infinitary sentences must not
be treated as one countable conjunction without a countability
condition. The precise formulation matters because there are
uncountably many infinitary sentences.
The first-order cardinal conjecture, its perfect-set
strengthening and the dated \(\omega\)-Vaught comparison
are distinguished in
\cref{def:p03-vaught-variants,def:p03-vaught-scott-comparison}.
The edition's selected status is in
\cref{sec:p03-vaught-status-2026}; no different rank convention
or later version is substituted here.

\subsection{Infinitary truth in the coding space}
\label{sec:p03-morley-borel-truth}

\begin{hbdefinition}{def:p03-morley-infinitary}{Countable connectives and formula rank}
The logic \(L_{\omega_1,\omega}\) permits countable
conjunctions and disjunctions, each with only finitely
many free variables in total, and finite strings of
individual quantifiers. Formula trees are well founded.
Define the auxiliary quantifier rank by
\[
 \operatorname{qr}(\theta)=0\ \text{for atomic }\theta,\quad
 \operatorname{qr}(\neg\varphi)=\operatorname{qr}(\varphi),
 \quad
 \operatorname{qr}\!\left(\bigwedge_i\varphi_i\right)
              =\sup_i\operatorname{qr}(\varphi_i),
 \quad
 \operatorname{qr}(\exists y\,\varphi)
              =\operatorname{qr}(\varphi)+1.
\]
Disjunction has the same rank clause as conjunction.
All these ranks are countable ordinals.
A set is Borel if it belongs to the sigma-algebra
generated by the open sets.
\end{hbdefinition}

\begin{hblemma}{thm:p03-morley-borel-truth}{proved=annex:p3.morley.borel-truth}{Each fixed infinitary formula defines a Borel truth set}
For a fixed \(L_{\omega_1,\omega}\)-formula
\(\varphi(\bar x)\) and each fixed tuple
\(\bar a\in\omega^n\),
\[
 \{M\in X_L:M\models\varphi(\bar a)\}
\]
is Borel. The satisfaction relation is also Borel in
\(X_L\times\omega^n\), with the tuple space discrete.
The same holds in each finite-domain component.
Every countable collection of infinitary sentences can
be conjoined into one such sentence, including the
empty collection with conjunction \(\top\).
Its fixed-domain model class is Borel.
\end{hblemma}

\subsection{Transfinite back-and-forth and the Scott construction}
\label{sec:p03-morley-scott}

\begin{hbdefinition}{def:p03-morley-approximants}{Pointed back-and-forth approximants}
For equal-length finite tuples in structures \(M,N\), define
\((M,\bar a)\equiv_\alpha(N,\bar b)\), for \(\alpha<\omega_1\):
\begin{enumerate}
\item At \(0\), the tuples agree on every atomic formula
      in their slots.
\item At \(\alpha+1\), they are \(\equiv_\alpha\)-equivalent,
      and every \(c\in M\) has some \(d\in N\) with
      \((M,\bar a c)\equiv_\alpha(N,\bar b d)\),
      while every \(d\in N\) has some such \(c\in M\).
\item At a nonzero limit \(\lambda\), all
      \(\equiv_\beta\), \(\beta<\lambda\), hold.
\end{enumerate}
Tuples may repeat; atomic equality identifies the same
repetitions on each side. Empty tuples are included.
Write \(E_\alpha(M,N)\) for the empty-tuple relation;
its Borel claims below concern codes on \(\omega\).
This auxiliary ordinal analysis is not an identification
with the parameterized or unparameterized Scott-rank
conventions in \cref{def:p03-vaught-scott-comparison}.
\end{hbdefinition}

\begin{hblemma}{thm:p03-morley-approximant-properties}{proved=annex:p3.morley.approximants}{Equivalence, Borelness and preservation of ranked formulas}
The pointed relations are equivalence relations in each
tuple arity and decrease with \(\alpha\).
For coded structures on \(\omega\), each fixed
countable-ordinal relation is Borel.
If \((M,\bar a)\equiv_\alpha(N,\bar b)\), the tuples
agree on every infinitary formula of quantifier rank
at most \(\alpha\).
\end{hblemma}

\begin{hbdefinition}{def:p03-morley-scott-formulas}{The formulas attached to one countable structure}
Fix an at most countable \(M\) and a finite tuple \(\bar a\).
Let \(\theta_{\bar a}^{0}(\bar x)\) be the conjunction
of the atomic formulas true at \(\bar a\) and the
negations of those false there. Define
\[
\begin{split}
 \theta_{\bar a}^{\alpha+1}(\bar x)
   ={}&\theta_{\bar a}^{\alpha}(\bar x)
       \land\bigwedge_{c\in M}\exists y\,
                         \theta_{\bar a c}^{\alpha}(\bar x,y)\\
      &\land\forall y\bigvee_{c\in M}
                         \theta_{\bar a c}^{\alpha}(\bar x,y),\\
 \theta_{\bar a}^{\lambda}(\bar x)
   ={}&\bigwedge_{\beta<\lambda}
                         \theta_{\bar a}^{\beta}(\bar x).
\end{split}
\]
All bound variables are renamed to avoid capture.
Each indexing set is countable and each formula has
only its displayed finite slots free.
\end{hbdefinition}

\begin{hblemma}{thm:p03-morley-scott-stabilization}{proved=annex:p03-morley-scott-stabilization}{Semantic representation and a countable stabilization stage}
For every structure \(N\),
\[
 N\models\theta_{\bar a}^{\alpha}(\bar b)
       \quad\Longleftrightarrow\quad
 (M,\bar a)\equiv_\alpha(N,\bar b).
\]
There is a countable ordinal \(\delta_M\) at which
all finite-tuple approximants \emph{within \(M\)}
stabilize:
\[
 (M,\bar a)\equiv_{\delta_M}(M,\bar b)
       \quad\Longrightarrow\quad
 (M,\bar a)\equiv_{\delta_M+1}(M,\bar b)
\]
for all equal-length tuples.
This countable stage depends on \(M\); no uniform
countable bound for all structures is asserted.
\end{hblemma}

\begin{hbtheorem}{thm:p03-morley-scott-sentence}{proved=annex:p3.morley.scott-sentence}{A full countable-isomorphism sentence}
The \(L_{\omega_1,\omega}\)-sentence
\[
 \sigma_M=
 \theta_\varnothing^{\delta_M}
 \land\bigwedge_{n<\omega}\ \bigwedge_{\bar a\in M^n}
   \forall\bar x\left(
       \theta_{\bar a}^{\delta_M}(\bar x)\to
       \theta_{\bar a}^{\delta_M+1}(\bar x)\right)
\]
is true in \(M\) and, among nonempty at most countable
structures, has exactly the models isomorphic to \(M\).
It can also be formed after naming any fixed finite
tuple, in which case it characterizes the pointed
isomorphism type, including repeated named coordinates.
It asserts uniqueness among countable models, not
among arbitrary uncountable models.
\end{hbtheorem}

\begin{hbcorollary}{thm:p03-morley-approximant-isomorphism}{proved=annex:p3.morley.approximant-isomorphism}{Isomorphism is the intersection of all countable approximants}
For at most countable \(M,N\),
\[
 M\cong N
       \quad\Longleftrightarrow\quad
 E_\alpha(M,N)\ \text{for every }\alpha<\omega_1.
\]
Each fixed \(E_\alpha\) is Borel on \(X_L\), but
the intersection here is indexed by all countable
ordinals, not by one countable list.
Closure of the Borel sets under countable intersections
does not make this an assertion that the whole
isomorphism relation is Borel.
\end{hbcorollary}

\subsection{Individual classes and a perfect family}
\label{sec:p03-morley-isomorphism-classes}

\begin{hbtheorem}{thm:p03-morley-isomorphism-complexity}{proved=annex:p3.morley.isomorphism-complexity}{Borel individual classes and analytic global isomorphism}
Each individual isomorphism class in \(X_L\) is Borel.
The global relation \(\cong\) on \(X_L\) is analytic:
it is the projection of a closed subset of
\(S_\infty\times X_L\times X_L\), where \(S_\infty\)
is the Polish group of permutations of \(\omega\).
Here analytic means projection of a Borel subset of
a product with a Polish space.
Borelness of each separate class does not by itself
prove Borelness of the global relation or a spectrum
dichotomy.
\end{hbtheorem}

\begin{hbtheorem}{thm:p03-morley-perfect-example}{proved=annex:p3.morley.perfect-example}{A complete perfect-set witness in an incomplete theory}
Let \(L=\{P_i:i<\omega\}\) consist of unary relations,
and let \(U\) say that every \(P_i\) is constant on
the domain:
\(\forall x\forall y(P_i(x)\leftrightarrow P_i(y))\).
For \(\eta\in2^\omega\), let \(M_\eta\) have domain
\(\omega\) and interpret \(P_i\) by all of \(\omega\)
when \(\eta(i)=1\), and by the empty set otherwise.
The map \(\eta\mapsto M_\eta\) is a homeomorphism
onto a nonempty perfect subset of \(X_L\) consisting
of pairwise nonisomorphic \(U\)-models.
It witnesses \(I(K_U)=2^{\aleph_0}\).
The theory \(U\) is incomplete: different models
decide \(\forall xP_i(x)\) differently.
The example is not a proof of the complete-theory
conjecture. The same binary patterns on any fixed
finite nonempty domain give continuum many finite
isomorphism types in this infinite language.
\end{hbtheorem}


\section{Order indiscernibles}
\label{sec:p03-indiscernibles}

\begin{hbdefinition}{def:p03-order-indiscernibles}{Order indiscernibles}
Let \(A\subseteq\dom\mathcal M\), let \((I,<)\) be a linear order, and let
\((a_i)_{i\in I}\) be a sequence of tuples of one fixed finite length in
\(\mathcal M\).  The sequence is \emph{order indiscernible over \(A\)} when,
for every \(n<\omega\), every two increasing \(n\)-tuples of indices
\(i_1<\cdots<i_n\) and \(j_1<\cdots<j_n\), and every
\(\Lang(A)\)-formula \(\varphi(z_1,\ldots,z_n)\), where each \(z_r\) is a
tuple of variables having the same length as \(a_i\),
\[
 \mathcal M\entails\varphi(a_{i_1},\ldots,a_{i_n})
 \quad\Longleftrightarrow\quad
 \mathcal M\entails\varphi(a_{j_1},\ldots,a_{j_n}).
\]
With \(A=\varnothing\) one simply says \emph{order indiscernible}.  This is a
property of an indexed sequence: unlike homogeneity, it compares increasing
finite subsequences by their order type rather than extending arbitrary
partial elementary maps.  The distinction is needed by the construction
theorem that follows.
\index{indiscernibles!order}
\end{hbdefinition}

\begin{hbtheorem}{thm:p03-ehrenfeucht-mostowski}{cited=EhrenfeuchtMostowski1956}{Ehrenfeucht--Mostowski indiscernibles}
Let \(T\) be a first-order theory with an infinite model, and let \((I,<)\)
be any linear order.  Some model of \(T\) contains a sequence of distinct
order indiscernibles indexed in order type \(I\).
\end{hbtheorem}

The cited construction combines compactness with finite Ramsey arguments.  At
this statement ceiling no Skolem expansion, coherent
blueprint, or claim that the model is generated by the indiscernibles is
being made \cite{EhrenfeuchtMostowski1956,Marker2002}.

To state the construction that lies behind applications, first take a
Skolem expansion \(T^{\mathrm{sk}}\) of \(T\).  An
\emph{Ehrenfeucht--Mostowski blueprint for linear orders} is a coherent
family \(\Phi=(p_n)_{1\leq n<\omega}\) of complete quantifier-free
\(n\)-types in the Skolem language: each \(p_n\) asserts that its variables
are distinct, and every increasing \(m\)-subtuple of a realization of
\(p_n\) realizes \(p_m\).  Coherence is needed because the same finite
suborder occurs inside many larger ones.  The blueprint is metalinguistic
data for building models, not an additional predicate or an object-language
theory; it differs from one indiscernible sequence by prescribing all finite
increasing types uniformly, and it is used to define the models in the next
statement.

For a linear order \(I\), \(\operatorname{EM}(I,\Phi)\) denotes the
\(\Lang\)-reduct of the Skolem hull of a sequence
\((a_i)_{i\in I}\) whose increasing tuples realize \(\Phi\).  This notation
is eliminable in favor of that hull construction; its only use here is to
state the functorial and spectrum consequences without repeatedly describing
the Skolem expansion.

\begin{hbtheorem}{thm:p03-em-models}{cited=EhrenfeuchtMostowski1956}{Ehrenfeucht--Mostowski models}
For every first-order theory \(T\) with an infinite model there is a blueprint
\(\Phi\) such that \(\operatorname{EM}(I,\Phi)\entails T\) for every linear
order \(I\), and every order embedding \(I\to J\) induces an elementary
embedding of the corresponding \(\Lang\)-reducts.
\end{hbtheorem}

\begin{hbtheorem}{thm:p03-unstable-spectrum}{cited=Shelah1971NonIsomorphicModels}{Unstable theories have the maximal uncountable spectrum}
If a complete theory \(T\) is unstable, then for every uncountable cardinal
\(\kappa>\abs T\) it has \(2^\kappa\) pairwise nonisomorphic models of
cardinality \(\kappa\).
\end{hbtheorem}

Finite Ramsey theory and compactness supply the coherent blueprint; Skolem
closure makes maps induced by order embeddings elementary.  For the spectrum
conclusion, an order-property formula lets a refined blueprint recover enough
of the index order from the resulting model that a family of
\(2^\kappa\) suitable orders yields nonisomorphic models.  The strict bound
\(\kappa>\abs T\) is load-bearing: the countable random graph is unstable but
\(\aleph_0\)-categorical, so no corresponding countable-cardinal claim is
valid \cite{EhrenfeuchtMostowski1956,Shelah1971NonIsomorphicModels,Hodges1993}.

\section{Stability, dividing, forking, and simplicity}
\label{sec:p03-stability}

For a parameter set \(A\) in a sufficiently saturated model of the complete
theory \(T\), let \(S_n(A)\) be the set of complete \(n\)-types over \(A\)
consistent with \(T\), and let
\[
\operatorname{tp}(a/A)
\defeq\{\varphi(x,\bar b):\bar b\text{ is from }A
              \text{ and }\varphi(a,\bar b)\text{ holds}\}.
\]
If \(b\) is a finite tuple, \(Ab\) abbreviates the union of \(A\) with its
entries.  For a complete type \(p\) over \(C\) and \(B\subseteq C\),
\(p\mathbin{\upharpoonright}B\) is the subset of formulas in \(p\) whose
parameters lie in \(B\).
These are parameter-set versions of the inherited type and type-space
definitions.  They are metalinguistic sets of formulas, not new sorts; unlike
\(S_n(T)\), they permit constants from \(A\).  The notation is required by
the stability-spectrum count and the nonforking laws below.

\begin{hbdefinition}{def:p03-stability}{Order property and stability}
Let \(T\) be complete.  A formula \(\varphi(x;y)\) has the \emph{order
property in \(T\)} when some \(\mathcal M\entails T\) contains tuples
\((a_i)_{i<\omega}\) and \((b_j)_{j<\omega}\) such that
\[
 \mathcal M\entails\varphi(a_i;b_j)
 \quad\Longleftrightarrow\quad i<j.
\]
The theory \(T\) is \emph{stable} when no formula has the order property in
\(T\).  Thus stability is a property of a complete theory defined through
all its formulas, not merely a claim that one selected model has few
automorphisms or that the theory has few models.  This exact obstruction is
used in the overview table below.
\index{order property}
\index{stability}
\end{hbdefinition}

\begin{hbdefinition}{def:p03-dividing-forking}{Dividing and forking}
Let \(T\) be complete, \(\mathcal M\entails T\), \(A\subseteq\dom\mathcal M\),
and let \(b\) be a finite tuple from \(\mathcal M\).
\begin{enumerate}
  \item A formula \(\varphi(x;b)\) \emph{divides over \(A\)} when, in some
        elementary extension of \(\mathcal M\), there is an
        \(A\)-indiscernible sequence \((b_i)_{i<\omega}\) with \(b_0=b\)
        for which \(\{\varphi(x;b_i):i<\omega\}\) is inconsistent.
  \item A formula \(\psi(x;c)\), with \(c\) from \(\mathcal M\),
        \emph{forks over \(A\)} when some elementary extension
        \(\mathcal N\succcurlyeq\mathcal M\) contains parameter tuples
        \(b_1,\ldots,b_m\), where \(m\geq1\), and formulas
        \(\varphi_j(x;b_j)\), each dividing
        over \(A\), such that
        \[
          \mathcal N\entails\forall x\,
          \bigl(\psi(x;c)\to\bigvee_{j=1}^{m}\varphi_j(x;b_j)\bigr).
        \]
  \item If \(A\subseteq B\subseteq\dom\mathcal M\), a complete type over
        \(B\) forks over \(A\) when it contains a formula that forks over
        \(A\); otherwise it is \emph{nonforking over \(A\)}.
\end{enumerate}
Dividing is witnessed by one indiscernible inconsistency pattern; forking is
its finite-disjunction closure.  The terms are relations among formulas or
types and parameter sets, not new kinds of models.  Their difference is
needed because they coincide only under additional hypotheses, notably in
simple theories.
\index{dividing}
\index{forking}
\index{nonforking}
\end{hbdefinition}

\begin{hbdefinition}{def:p03-simple-theory}{Simple theory}
Let \(T\) be complete.  Write
\(\abs T\) with the chapter-wide convention fixed above.
In the local-character presentation used here, \(T\) is \emph{simple} when,
for every \(\mathcal M\entails T\), every parameter set
\(B\subseteq\dom\mathcal M\), and every complete type \(p(x)\) over \(B\)
consistent with \(\mathcal M\), some \(A\subseteq B\) with
\(\abs A\leq\abs T+\abs x\) satisfies: \(p\) does not fork over \(A\).
For a finite tuple \(x\), this is the usual bound
\(\abs A\leq\abs T\).  Simplicity therefore names local character of the
forking relation; it is weaker than stability, as the example below records,
and is used to locate forking in the table.  The cited sources establish the
equivalence of this presentation with their standard definitions; no stronger
simple-theory machinery is asserted here.
\index{simple theory}
\end{hbdefinition}

Here \emph{the random graph} means the complete theory of the countable simple
graph in which, for every two disjoint finite vertex sets \(U,V\), some vertex
outside \(U\cup V\) is adjacent to every member of \(U\) and to no member of
\(V\).  This distinguishes the theory used here from probability distributions
on finite graphs.
Every stable theory is simple, and the cited simple-theory development gives
symmetry and transitivity for nonforking.  The random graph is simple but
unstable, so instability alone does not imply nonsimplicity
\cite{Shelah1971Stability,TentZiegler2012,Kim1998}.
\index{random graph}

\subsection{The stability spectrum and nonforking calculus}

\begin{hbdefinition}{def:p03-stability-spectrum}{Cardinal stability, \(\omega\)-stability, and superstability}
Let \(T\) be complete and work in a sufficiently saturated ambient model.
For any infinite cardinal \(\kappa\), the theory is
\emph{\(\kappa\)-stable} when
\[
  \abs{S_1(A)}\leq\kappa
  \qquad\text{for every parameter set }A\text{ with }\abs A\leq\kappa.
\]
It is \emph{\(\omega\)-stable} when it is \(\aleph_0\)-stable.  A stable
theory is \emph{superstable} when every complete finite-tuple type over a set
\(B\) does not fork over some finite subset of \(B\).

The cardinal spectrum is normally analyzed at \(\kappa\geq\abs T\), but the
type-counting predicate itself is defined for every infinite \(\kappa\); this
makes \(\omega\)-stability meaningful exactly as written, including before a
countability hypothesis is imposed on the language.

Cardinal stability measures the size of type spaces, while superstability
states finite local character for the forking relation.  These are properties
of a complete theory, not saturation properties of one model.  They are not
redundant with the order-property definition of stability: the latter says
whether the stability spectrum is nonempty, whereas these conditions locate
the theory within that spectrum.  The distinctions are used in the hierarchy
and rank theorems below.
\index{stability!cardinal spectrum}\index{omega-stability@\(\omega\)-stability}
\index{superstability}
\end{hbdefinition}

\begin{hbtheorem}{thm:p03-nonforking-calculus}{cited=TentZiegler2012}{Nonforking in stable theories}
Let \(T\) be stable.
\begin{enumerate}
  \item If \(A\subseteq B\) and \(p\) is a complete finite-tuple type over
        \(A\), then \(p\) has an extension to \(B\) that does not fork over
        \(A\).
  \item For tuples \(a,b\), the type \(\operatorname{tp}(a/Ab)\) does not
        fork over \(A\) iff \(\operatorname{tp}(b/Aa)\) does not fork over
        \(A\).
  \item If \(A\subseteq B\subseteq C\) and \(p\) is a complete type over
        \(C\), then \(p\) does not fork over \(A\) iff both \(p\) does not
        fork over \(B\) and \(p\mathbin{\upharpoonright}B\) does not fork
        over \(A\).
\end{enumerate}
Thus nonforking has extension, symmetry, and transitivity.  Moreover,
\(\omega\)-stable theories are superstable and superstable theories are
stable; neither converse holds.
\end{hbtheorem}

The theorem makes nonforking usable as an independence relation rather than
leaving it at the negative definition ``does not contain a forking formula.''
The stability hypothesis is essential for this package as stated; the simple-
theory generalization requires its own forking development
\cite{Shelah1971Stability,TentZiegler2012}.

\section{Morley rank and uncountable categoricity}
\label{sec:p03-morley-rank}

\begin{hbdefinition}{def:p03-morley-rank}{Morley rank and total transcendence}
Let \(T\) be complete, and let \(X\) be a definable subset of a sufficiently
saturated model of \(T\), with parameters allowed.  Define the relation
\(\operatorname{RM}(X)\geq\alpha\) by transfinite recursion:
\begin{enumerate}
  \item \(\operatorname{RM}(X)\geq0\) iff \(X\ne\varnothing\);
  \item \(\operatorname{RM}(X)\geq\alpha+1\) iff \(X\) contains infinitely
        many pairwise disjoint definable subsets \(X_i\) with
        \(\operatorname{RM}(X_i)\geq\alpha\); and
  \item at a limit \(\lambda\), \(\operatorname{RM}(X)\geq\lambda\) iff it
        is at least every \(\alpha<\lambda\).
\end{enumerate}
Set \(\operatorname{RM}(\varnothing)=-\infty\).  For nonempty \(X\), the
\emph{Morley rank} \(\operatorname{RM}(X)\) is the least \(\alpha\) for which
it is not at least \(\alpha+1\), and is \(\infty\) if there is no such ordinal.
A theory is \emph{totally transcendental} when every definable set has ordinal
Morley rank (with the empty-set convention just stated).

The rank is needed to measure repeated definable splitting, which cardinal
counts of complete types do not record locally.  It is an ordinal-valued
metalinguistic invariant of definable sets; total transcendence is a property
of the theory.  Because parameters are allowed,
\(\operatorname{RM}(X)=0\) exactly for nonempty finite \(X\); higher ranks
measure definable splitting rather than cardinality, and rank \(\infty\) is
not a new infinite cardinal.  The distinction is used in the categoricity
outline below.
\index{Morley rank}\index{total transcendence}
\end{hbdefinition}

\begin{hbdefinition}{def:p03-morley-geometry}{Algebraic closure, Morley degree, and dimension}
For a parameter set \(A\), the model-theoretic algebraic closure is
\[
 \operatorname{acl}(A)
 \defeq\{b:b\text{ belongs to some finite set definable over }A\}.
\]
If \(X\) has ordinal Morley rank \(\alpha\), its \emph{Morley degree} is the
largest \(d<\omega\) for which \(X\) has \(d\) pairwise disjoint definable
subsets of rank \(\alpha\).  The maximum is finite: otherwise compactness
would give infinitely many such subsets and force rank at least
\(\alpha+1\).

A \emph{pregeometry} on a set \(D\) is a closure operator
\(\operatorname{cl}:\mathcal P(D)\to\mathcal P(D)\) that is extensive,
monotone, idempotent, has finite character, and satisfies exchange:
\[
 a\in\operatorname{cl}(X\cup\{b\})\setminus\operatorname{cl}(X)
 \quad\Longrightarrow\quad
 b\in\operatorname{cl}(X\cup\{a\}).
\]
A basis is a maximal set whose members are not in the closure of the others;
its cardinality is the \emph{dimension}.  On a strongly minimal set \(D\)
defined over parameters \(A_0\), the relevant closure is
\(\operatorname{cl}(X)=\operatorname{acl}(A_0\cup X)\cap D\).
Finally, a model \(M\) is \emph{prime over} a parameter set \(B\subseteq M\)
when every model of \(T\) containing \(B\) receives an elementary embedding
of \(M\) fixing \(B\) pointwise.

These notions are needed to state what the categoricity outline counts:
algebraic closure records finite definability, degree counts top-rank pieces,
and pregeometry dimension counts independent generators.  They stand,
respectively, as a closure operation, a finite invariant of a definable set,
and a cardinal invariant of a closure space; none is a synonym for ordinary
cardinality or Morley rank.  Prime-over-a-base compares models rather than
points of the pregeometry.  All four distinctions are used in the
Baldwin--Lachlan analysis below.
\index{algebraic closure!model-theoretic}\index{Morley degree}
\index{pregeometry}\index{dimension!strongly minimal set}
\end{hbdefinition}

\begin{hbdefinition}{def:p03-strongly-minimal}{Strongly minimal definable set}
An infinite definable set \(D\) is \emph{strongly minimal} when, in every
elementary extension, every subset of \(D\) definable with parameters is
finite or cofinite in \(D\).  The extension clause is needed because a
coincidental finite/cofinite pattern in one model is not preserved
automatically.  Strong minimality is a property of a definable set, not of an
entire theory; it differs from minimality in cardinality and from
\(\aleph_0\)-categoricity.  Algebraic closure on \(D\) then yields the
pregeometry whose dimension is used by Baldwin--Lachlan.
\index{strongly minimal set}
\end{hbdefinition}

\begin{hbtheorem}{thm:p03-total-transcendence}{cited=Marker2002}{Rank and \(\omega\)-stability}
For a complete theory in a countable language, total transcendence is
equivalent to \(\omega\)-stability.
\end{hbtheorem}

\begin{hbtheorem}{thm:p03-baldwin-lachlan}{cited=BaldwinLachlan1971}{Baldwin--Lachlan refinement}
Let \(T\) be a complete theory in a countable language that is categorical in
one uncountable cardinal.  Then \(T\) is \(\omega\)-stable; after naming
finitely many parameters, a strongly minimal definable set controls its
uncountable models by pregeometry dimension.  Consequently \(T\) is
categorical in every uncountable cardinal, and
\[
  I(T,\aleph_0)=1
  \qquad\text{or}\qquad
  I(T,\aleph_0)=\aleph_0.
\]
\end{hbtheorem}

The proof architecture behind Morley's categoricity theorem is now visible.
Uncountable categoricity first rules out excessive type splitting and yields
\(\omega\)-stability, equivalently ordinal Morley rank.  A definable set of
minimal positive rank and degree becomes strongly minimal after finitely many
parameters are named.  Algebraic closure on that set is a pregeometry; bases
have well-defined dimension, and the prime-over-a-basis analysis shows that
uncountable models of the same cardinal have the same dimension and hence are
isomorphic.  Baldwin and Lachlan sharpen this analysis to the displayed
countable-model dichotomy.  The hard existence and coordinatization lemmas are
those of the cited proofs; this paragraph is a proof outline, not a replacement
for them \cite{Morley1965,BaldwinLachlan1971,Marker2002}.

\section{Morley's Original Categoricity-in-Power Argument}
\label{sec:p03-n-a-006-01}

Morley's 1965 argument proves more than transfer of categoricity: if a
countable theory is categorical in one uncountable cardinal, every
uncountable model is saturated. Its machinery consists of a derivative
of all parameter type spaces together, a Ramsey construction of models
with few types realized over countable sets, and prime models arranged
compatibly along chains. These are the original five sections of
\emph{Categoricity in Power}, including its three rank examples and
final rank bound \cite[pp.~514--538]{Morley1965}. Strongly minimal
pregeometry and prime-over-basis coordinatization are later developments;
they are not substituted for this proof.

\subsection{Normalization and the category of parameter systems}
\label{sec:p03-morley-normalization}

\begin{hbconvention}{def:p03-morley-standing}{The original standing assumptions}
Let \(\Sigma\) be a satisfiable complete theory in an at most countable
first-order language, with an infinite model. First replace any
functions and constants by their total single-valued
graph relations and singleton relations, translating term formulas.
In the displays below \(\Sigma\) denotes this relational presentation.
For every relational formula
\(\varphi(v_0,\ldots,v_{n-1})\), introduce a relation symbol whose
interpretation is its truth set, retaining equality. The resulting
definitional relational expansion is \(\Sigma^*\). In this presentation
\(T=\Sigma^*\), and \(\mathcal N(T)\) is the class of induced
subsystems of models of \(T\). Subsystems need not themselves be models.
The empty parameter system is permitted for type spaces.

A \emph{monomorphism} preserves and reflects every relation and is
injective. In the expanded language it is precisely a partial elementary
map of the original structures. Thus unions and amalgams below preserve
the original elementary information, not merely atomic information in
the unexpanded language. As in the source, \(|A|\) denotes the
universe of the system \(A\), whereas \(\kappa(X)\) denotes the
cardinality of a set \(X\), and \(\kappa(A)=\kappa(|A|)\) is the
cardinality, or power, of the system. A bare \(\kappa\) denotes a
cardinal; the argument distinguishes the cardinality operation
\(\kappa(\cdot)\). The source's equality of cardinals modulo
\(\aleph_0\) means equality after adding \(\aleph_0\).
The ambient set theory includes Choice.
\end{hbconvention}

\begin{hbtheorem}{thm:p03-morley-normalization}{proved=annex:p3.morley.normalization}{Definitional expansion, chains, and amalgamation}
Every \(\Sigma\)-model has exactly one expansion to a \(T\)-model.
Isomorphisms, partial elementary maps, model cardinalities, saturation,
and categoricity transfer in both directions. Every expanded formula is
equivalent modulo \(T\) to an atomic relation formula.
The theory \(T\) is complete. The class \(\mathcal N(T)\) is closed
under increasing unions, has joint embedding, and has amalgamation over
every common subsystem. Increasing unions of \(T\)-models are models.
\end{hbtheorem}

If an initial theory is categorical in an infinite cardinal and has no
finite models, completeness follows: two completions would, by
L\"owenheim--Skolem, give non-elementarily-equivalent models in that
cardinal. If finite models are allowed, the conclusion applies to its
infinite-model part, obtained by adding the sentences asserting at least
\(n\) elements for each finite \(n\). Finite models must not be folded
into the source's standing infinite complete theory.

\begin{hbdefinition}{def:p03-morley-type-space}{Formula algebras, points, and restriction maps}
For \(A\in\mathcal N(T)\), let \(T(A)\) be \(T\) together with the
diagram naming every element of \(A\). Because the language is
Morleyized, this is a complete parameter theory. Let \(F(A)\) be the
Boolean algebra of formulas with at most the free variable \(v_0\),
modulo equivalence in \(T(A)\). Its Stone space \(S(A)\) consists of
complete one-element types over \(A\), with clopen basis
\[
 U_\varphi=\{p\in S(A):\varphi\in p\}.
\]
A point is realized by \(b\) in an extension when every formula in it
holds of \(b\). For a monomorphism \(f:A\to B\), substitute
\(f(a)\) for named parameters to obtain the Boolean embedding
\(F(A)\to F(B)\). Its dual is the restriction map
\(f^*:S(B)\to S(A)\). For an inclusion we write \(i_{AB}^*\).
These maps, with all the spaces \(S(A)\), constitute the source's
category of continuous onto maps. Inverse images always refer to
actual restriction maps, not arbitrary continuous maps.
\end{hbdefinition}

\begin{hblemma}{thm:p03-morley-types}{proved=annex:p3.morley.types}{Realization and topology of parameter types}
Every \(S(A)\) is compact Hausdorff and zero-dimensional, with a
clopen basis of size at most \(\kappa(A)+\aleph_0\).
Every point is realized in an extension, and one extension realizes
all points at once. All restriction maps are continuous onto maps.
Two elements in possibly different extensions have the same type over
\(A\) exactly when the map identifying them and fixing \(A\) is a
monomorphism on the induced systems. Each basic clopen is the inverse
image of a basic clopen over a finite parameter subsystem.
\end{hblemma}

\subsection{The global derivative, rank, and finite degree}
\label{sec:p03-morley-global-derivative}

\begin{hbdefinition}{def:p03-morley-global-rank}{Morley's derivative across all extensions}
Simultaneously for all \(A\in\mathcal N(T)\), put
\[
 S^\alpha(A)=S(A)\setminus\bigcup_{\beta<\alpha}
                  \operatorname{Tr}^\beta(A).
\]
A point \(p\in S^\alpha(A)\) belongs to
\(\operatorname{Tr}^\alpha(A)\) when, for every monomorphism
\(f:A\to B\), every point of
\[
       (f^*)^{-1}(p)\cap S^\alpha(B)
\]
is isolated in \(S^\alpha(B)\). The rank of such a point is
\(\alpha\); points of \(\operatorname{Tr}^0(A)\) are called
\emph{algebraic}. A point surviving every stage has rank \(\infty\).
An isolated point of \(S(A)\) need not be algebraic: isolation in the
single space is weaker than isolation of every extension fiber.
The identity map is included among the tested maps.
\end{hbdefinition}

\begin{hbtheorem}{thm:p03-morley-derivative}{proved=annex:p3.morley.derivative}{Closedness, exact image, and finite localization}
For every ordinal \(\alpha\), \(S^\alpha(A)\) is closed and compact.
For every monomorphism \(f:A\to B\),
\[
 f^*(S^\alpha(B))=S^\alpha(A).
\]
For \(p\in S^\alpha(A)\),
\[
 p\in\operatorname{Tr}^\alpha(A)
 \quad\Longleftrightarrow\quad
 (f^*)^{-1}(p)\cap S^\alpha(B)
           \subseteq\operatorname{Tr}^\alpha(B).
\]
If \(p\) has ordinal rank \(\alpha\), some finite \(E\subseteq A\)
has \(i_{EA}^*(p)\) of exactly that rank.
\end{hbtheorem}

\begin{hbtheorem}{thm:p03-morley-degree}{proved=annex:p3.morley.degree}{Uniform finite multiplicity and its additive law}
If \(p\) has rank \(\alpha\), there is a positive integer bounding
\(\kappa((f^*)^{-1}(p)\cap S^\alpha(B))\) for every
\(f:A\to B\). The maximum of these integers is its degree
\(d(p)\). For every such map,
\[
 d(p)=\sum_{q\in(f^*)^{-1}(p)\cap S^\alpha(B)}d(q).
\]
In particular, an extension of the same rank has no larger degree;
if it has the same degree, it is the only extension of that rank.
Degree is a finite multiplicity, not the ordinal rank or the
cardinality of a basis.
\end{hbtheorem}

\begin{hbdefinition}{def:p03-morley-total-transcendence}{Stabilization and total transcendence}
Let \(\alpha_T\) be the least ordinal at which every global derivative
has stabilized: \(S^{\alpha_T}(A)=S^\beta(A)\) for every
\(\beta\geq\alpha_T\) and every \(A\). The theory is
\emph{totally transcendental} when this common residual space is empty.
In that case every type over every subsystem has ordinal rank.
\end{hbdefinition}

\begin{hbtheorem}{thm:p03-morley-stabilization}{proved=annex:p3.morley.stabilization}{Uniform stabilization and the countable rank bound}
The ordinal \(\alpha_T\) exists and is smaller than
\((2^{\aleph_0})^+\). If \(S^\alpha(A)\) is empty for one
subsystem, it is empty for every subsystem. The first empty stage is a
successor ordinal. If \(T\) is totally transcendental, then
\(\alpha_T<\omega_1\).
\end{hbtheorem}

\begin{hbtheorem}{thm:p03-morley-countability}{proved=annex:p3.morley.countability}{The original countability characterization}
If \(T\) is totally transcendental,
\[
 \kappa(S(A))\leq\kappa(A)+\aleph_0,
 \qquad
 \kappa(S(A))+\aleph_0=\kappa(A)+\aleph_0.
\]
Conversely, \(T\) is totally transcendental if \(S(A)\) is countable
for every countable \(A\in\mathcal N(T)\). If it is not totally
transcendental, some countable subsystem has \(2^{\aleph_0}\) types.
For the original countable theory this is equivalently
\(\omega\)-stability: countably many complete one-types over every
countable parameter set in a model.
\end{hbtheorem}

\begin{hbdefinition}{def:p03-morley-definable-rank}{Definable-set notation and its scope}
For a definable set \(X\) in finitely many variables, with arbitrary
parameters allowed, define \(\operatorname{RM}(X)\geq0\) when
\(X\ne\varnothing\); define
\(\operatorname{RM}(X)\geq\alpha+1\) when some elementary extension
has infinitely many pairwise disjoint definable subsets of \(X\),
each of rank at least \(\alpha\); and at a limit require every
earlier inequality. The empty set has rank \(-\infty\); a nonempty
set has its exact ordinal rank when the next inequality fails,
and rank \(\infty\) if none fails. Equivalently one works in a
sufficiently saturated ambient extension and allows its parameters.
For ordinal rank \(\alpha\), its Morley degree is the largest
number of disjoint definable subsets of rank \(\alpha\).
\end{hbdefinition}

\begin{hbtheorem}{thm:p03-morley-rank-bridge}{proved=annex:p3.morley.rank-bridge}{The bridge to definable-set Morley rank}
In a totally transcendental theory, the rank of a one-variable
\(\varphi\)-set over \(A\) is the largest rank of a point of
\(U_\varphi\subseteq S(A)\). Its degree is the sum of the degrees
of its points of that largest rank. These are exactly the splitting
rank and finite degree just defined. The same construction applies
to every finite tuple space. For countable theories, all definable
sets have ordinal Morley rank exactly when the theory is
\(\omega\)-stable. A rank-zero definable set is nonempty finite;
its degree is its number of elements. Rank and degree are not
algebraic-closure dimension.
\end{hbtheorem}

\subsection{The original three examples}
\label{sec:p03-morley-original-examples}

\begin{hbexample}{def:p03-morley-field-data}{Algebraically closed fields of a fixed characteristic}
Let \(\Sigma=\mathrm{ACF}_p\), where \(p=0\) or a fixed prime, and
let \(K(A)\) be the field generated by a parameter subsystem \(A\).
Over \(A\), the algebraic one-types correspond to irreducible
polynomials over \(K(A)\); there is one further type, asserting
nonvanishing of every nonzero polynomial over \(K(A)\).
This last type is the transcendental type. Morley's original example
uses characteristic zero; the same calculation works for each fixed
positive characteristic.
\end{hbexample}

\begin{hbtheorem}{thm:p03-morley-field-example}{proved=annex:p3.morley.field-example}{Field ranks, degrees, and the cardinal distinction}
The polynomial type for an irreducible polynomial has rank zero
and degree equal to its number of distinct roots (its separable
degree); in characteristic zero this is its polynomial degree.
The transcendental type has rank
one and degree one. Thus \(S^1(A)\) consists of that single type,
\(S^2(A)=\varnothing\), and \(\alpha_T=2\).
Every uncountable algebraically closed field of fixed characteristic
has transcendence degree equal to its cardinality; hence these
fields are categorical in every uncountable cardinal. Their
countably infinite models have the distinct transcendence degrees
\(0,1,2,\ldots,\aleph_0\), so they are not
\(\aleph_0\)-categorical. This numerical dimension is not the
one-variable rank one.
\end{hbtheorem}

\begin{hbexample}{def:p03-morley-tuple-code-data}{A categorical theory with larger finite rank}
Fix \(n\geq1\). There is an infinite unary part \(R_0\), and every
element outside it codes exactly one ordered \(n\)-tuple of distinct
elements of \(R_0\). Conversely each such tuple has exactly one
code, recorded by a relation \(R_1\) of arity \(n+1\).
For a parameter system \(B\), close it under taking code coordinates
and under forming codes; call the resulting set \(\overline B\).
\end{hbexample}

\begin{hbtheorem}{thm:p03-morley-tuple-code-example}{proved=annex:p3.morley.tuple-code-example}{Ranks in the tuple-code example}
This theory is categorical in every infinite cardinal. Elements of
\(\overline B\) have algebraic types. There is one type for a fresh
\(R_0\)-element, of rank one. A code type with exactly \(m\) of
its coordinates in \(\overline B\) has rank \(n-m\), for
\(m<n\). Consequently \(\alpha_T=n+1\). The calculation uses
extensions naming additional coordinates; counting isolated points
only in the initial space would miss these ranks.
\end{hbtheorem}

\begin{hbexample}{def:p03-morley-unary-data}{Closed subsets of Cantor space}
For a nonempty closed \(Y\subseteq2^\omega\), use unary predicates
\(R_j\) for \(j<\omega\). A finite Boolean predicate cell is empty
if its corresponding cylinder misses \(Y\), and infinite if that
cylinder meets \(Y\). These axioms give a complete theory: finite
formulas distinguish only finitely many predicate cells and equality
patterns. Write \(Y^{(\beta)}\) for the ordinary Cantor--Bendixson
derivatives, and \(\alpha_Y\) for their first stabilization stage.
\end{hbexample}

\begin{hbtheorem}{thm:p03-morley-unary-example}{proved=annex:p3.morley.unary-example}{The ordinary derivative inside the global derivative}
The named elements of \(A\) give the algebraic points of \(S(A)\).
The remaining points, one fresh-element type for each predicate
pattern in \(Y\), form \(S^1(A)\cong Y\). Restriction maps induce
homeomorphisms between these residual spaces, and
\[
 S^{1+\beta}(A)\cong Y^{(\beta)},\qquad
 \alpha_T=1+\alpha_Y.
\]
Here \(1+\beta\) is ordinal addition. The theory is totally
transcendental exactly when \(Y\) is countable, equivalently when
its perfect kernel is empty.
\end{hbtheorem}

\subsection{Ramsey models and total transcendence from categoricity}
\label{sec:p03-morley-ramsey-models}

\begin{hbdefinition}{def:p03-morley-skolem-hull}{Ordered indiscernibles and generated hulls}
An ordered set \((X,<)\) is a set of order indiscernibles in a
structure when increasing tuples of the same finite length satisfy
the same formulas. A countable Skolem expansion \(\Sigma^\#\)
is a conservative expansion by functions with universal axioms:
its reducts are \(\Sigma\)-models and every \(\Sigma\)-model
has an expansion. Its hull \(M(X,A)\) consists of all values of
terms on finite tuples from \(X\), including constant terms.
For elementary rather than merely relational preservation the
construction first Morleyizes and then Skolemizes, using
\(\Sigma^{*\#}\).
\end{hbdefinition}

\begin{hbtheorem}{thm:p03-morley-ramsey-em}{proved=annex:p3.morley.ramsey-em}{Ramsey, ordered indiscernibles, and the term-hull construction}
Every finite coloring of the \(n\)-element subsets of an infinite
set has an infinite homogeneous subset. Consequently, for any
theory with an infinite model and any ordered set \(X\), there is
a model containing \(X\) as order indiscernibles. In a Skolem
expansion, an order-preserving injection of the generators induces
a well-defined injective map of their term hull; a bijection induces
an automorphism. With the preliminary Morleyization these maps are
elementary on the original reduct. The hull has cardinality
\(\kappa(X)\) when \(X\) is infinite and the original
language is countable.
\end{hbtheorem}

\begin{hbtheorem}{thm:p03-morley-thin-models}{proved=annex:p3.morley.thin-models}{Thin models and the consequence of uncountable categoricity}
For every infinite cardinal \(\kappa\), a countable theory with
an infinite model has a model of size \(\kappa\) realizing only
countably many one-types over each countable subset of itself.
If the theory is categorical in some \(\kappa>\aleph_0\), it is
totally transcendental, hence \(\omega\)-stable.
Countable categoricity alone does not imply this: dense linear
order without endpoints has continuum many cut types over its
countable rational model.
\end{hbtheorem}

\begin{hbdefinition}{def:p03-morley-connected}{Connected and antisymmetric relations}
An \(n\)-ary relation is \emph{connected} on \(X\) if each
tuple of \(n\) distinct elements has a permutation satisfying it.
It is \emph{antisymmetric} in Morley's sense if each such tuple
also has a permutation not satisfying it. The second condition
is not the usual binary order axiom called antisymmetry.
\end{hbdefinition}

\begin{hbtheorem}{thm:p03-morley-no-order}{proved=annex:p3.morley.no-order}{No infinite permutation-sensitive relation}
In a totally transcendental theory no relation is both connected
and antisymmetric on an infinite subset of a model. The same holds
for a definable relation with a finite tuple of parameters.
\end{hbtheorem}

\subsection{Prime models compatible with chains}
\label{sec:p03-morley-prime-chains}

\begin{hbdefinition}{def:p03-morley-prime}{Prime extensions and unordered indiscernibles}
A model \(B\supseteq A\) is \emph{prime over} \(A\) when every
monomorphism of \(A\) into any model extends to a monomorphism of
\(B\) into that model. An \emph{indiscernible set over} \(A\)
is a set \(X\) disjoint from \(A\) for which every injective map
of \(A\cup X\) into itself fixing \(A\) is a monomorphism.
There is no distinguished order in this second notion.
\end{hbdefinition}

\begin{hbtheorem}{thm:p03-morley-prime-models}{proved=annex:p3.morley.prime-models}{Density and existence of prime extensions}
A subsystem \(A\) is a model exactly when its realized points
are dense in \(S(A)\). In a totally transcendental theory the
isolated points of each \(S(A)\) are dense: a minimum-rank point
of a nonempty open set is isolated in the whole space.
More generally, if these isolated points are dense for every
subsystem, every subsystem has a prime model over it, of size at
most \(\kappa(A)+\aleph_0\).
\end{hbtheorem}

\begin{hblemma}{thm:p03-morley-chain-types}{proved=annex:p3.morley.chain-types}{Coherent types and isolated lifts at limits}
Let \((A_\xi)_{\xi<\gamma}\) be an increasing nonempty chain,
with union \(A\), in a totally transcendental theory. A coherent
family \(p_\xi\in S(A_\xi)\) has eventually constant rank and
degree; it has a unique union type in \(S(A)\), of exactly those
eventual values. If \(p_0\) is isolated, it has coherent lifts
\(p_\xi\) that are isolated at every stage. One can include the
union stage; its lift is isolated there too.
\end{hblemma}

\begin{hbtheorem}{thm:p03-morley-chain-primes}{proved=annex:p3.morley.chain-primes}{Continuous chains of prime models}
For every increasing chain \((A_\xi)_{\xi<\gamma}\) in
\(\mathcal N(T)\) that is continuous at limit indices, there are
increasing models \((B_\xi)_{\xi<\gamma}\) with \(B_\xi\)
prime over \(A_\xi\), and
\(B_\delta=\bigcup_{\xi<\delta}B_\xi\) at limits.
The construction keeps each newly adjoined element isolated over
every later base together with the elements already adjoined.
This last condition is the transfer mechanism; separately chosen
prime models without it do not prove the theorem.
\end{hbtheorem}

\begin{hbtheorem}{thm:p03-morley-indiscernible-extraction}{proved=annex:p3.morley.indiscernible-extraction}{Large indiscernible sets inside the given system}
Suppose \(A\subseteq B\) are subsystems and
\(\kappa(A)<\kappa(B)=\kappa>\aleph_0\).
If \(\kappa\) is regular, \(B\setminus A\) contains an
indiscernible set over \(A\) of size \(\kappa\).
For singular \(\kappa\), for every \(\lambda<\kappa\) it
contains such a set of size greater than \(\lambda\).
This result assumes total transcendence and extracts its elements
inside \(B\), not merely inside an unspecified extension.
\end{hbtheorem}

\subsection{Saturation, omission transfer, and the final theorem}
\label{sec:p03-morley-saturation-transfer}

\begin{hbdefinition}{def:p03-morley-saturation}{Full and countable-subsystem saturation}
An infinite \(B\in\mathcal N(T)\) is saturated when every type
over each \(A\subseteq B\) with \(\kappa(A)<\kappa(B)\)
is realized in \(B\). For uncountable \(B\), saturation over
countable subsystems requires this only for countable \(A\).
Either property implies that \(B\) is a model. These are
realization properties, not an assertion that \(B\) realizes
every type over itself.
\end{hbdefinition}

\begin{hbtheorem}{thm:p03-morley-saturation-existence}{proved=annex:p3.morley.saturation-existence}{Uniqueness and the countable-saturation comparison model}
Two saturated models of the same infinite cardinality are
isomorphic. If \(T\) is totally transcendental and
\(\kappa>\aleph_0\), it has a model of cardinality \(\kappa\)
saturated over countable subsystems, with no assumption
\(\kappa^{\aleph_0}=\kappa\).
\end{hbtheorem}

\begin{hblemma}{thm:p03-morley-countable-omission-witness}{proved=annex:p3.morley.countable-omission-witness}{A countable witness with indiscernibles}
If an uncountable model \(B\) of a totally transcendental theory
is not saturated, there are a countable elementary submodel
\(A\subseteq B\), a subsystem \(A'\subseteq A\), an infinite
indiscernible set \(Y\subseteq A\setminus A'\) over \(A'\),
and a type \(q\in S(A')\) omitted in \(A\).
\end{hblemma}

\begin{hbtheorem}{thm:p03-morley-omission-transfer}{proved=annex:p3.morley.omission-transfer}{Transfer of a saturation failure to every uncountable size}
If a totally transcendental theory has any uncountable nonsaturated
model, then for each \(\kappa>\aleph_0\) it has a model of
cardinality \(\kappa\) that omits a type over a countable
subsystem. In particular that model is not saturated over
countable subsystems.
\end{hbtheorem}

\begin{hbtheorem}{thm:p03-morley-original-categoricity}{proved=annex:p3.morley.original-categoricity}{Morley's categoricity theorem and its stronger conclusion}
For a complete theory in a countable language, with an infinite
model, categoricity in one uncountable cardinal implies that
every uncountable model is saturated. Consequently it is
categorical in every uncountable cardinal. No continuum
hypothesis is required.
\end{hbtheorem}

The two comparison models in the categorical cardinal are an actual
countable-subsystem-saturated model and an actual model omitting a
countable-parameter type. Full saturation is concluded only after
the omission-transfer theorem rules out all uncountable failures.
This prevents replacing the hard transfer argument by the much
weaker observation that a suitably saturated extension exists.

\subsection{The original closing questions and the boundary of this source}
\label{sec:p03-morley-closing-questions}

The final page records questions rather than additional theorems
\cite[pp.~537--538]{Morley1965}. Historically, Morley asks whether an
uncountably categorical but not countably categorical theory has
exactly countably many countable models; whether such a theory can
be finitely axiomatizable; whether a language of size at most
\(\kappa>\aleph_0\) has categoricity transfer above \(\kappa\);
and whether a theory all of whose models have size at least
\(\kappa\) can be categorical in \(\kappa\).
For totally transcendental theories he asks for a model generated,
in the relevant symmetry sense, by a set all of whose permutations
extend to automorphisms; his footnote reports Silver's affirmative
answer. He also asks whether \(\alpha_T\) can ever be uncountable
(the proved total-transcendence case is negative), and what
conditions force it to be finite. He suggests categoricity and
finite axiomatizability as candidates, not as proved conditions.

These are the source's historical questions, not a declaration
that their 1965 status remains current. The later Baldwin--Lachlan
analysis and its countable-model consequences belong to their
own source presentation. The present proof uses neither an
unsupplied strongly minimal set nor an unsupplied prime-over-basis
theorem. The bibliographic works listed in Morley's references
retain their cited roles; that list does not allocate their whole
contents to this article's reconstruction.

\section{Strongly minimal geometry and the Baldwin--Lachlan refinement}
\label{sec:p03-n-a-006-02}

The refinement of uncountable categoricity by Baldwin and Lachlan uses
strongly minimal sets rather than identifying cardinality with an ordinal
rank. Their 1971 article treats the strongly minimal approach to categoricity
and then proves the countable-model dichotomy
\cite{BaldwinLachlan1971}. The geometric results below have explicit local
proofs. Their hypotheses do not include the assertion that every
uncountably categorical theory has already been coordinatized by such a
set: that assertion is a further theorem, not a definition of dimension.

\begin{hbconvention}{def:p03-bl-ambient}{Parameters and algebraicity}
Fix a complete first-order theory \(T\), a model \(\mathcal U\), and a finite
tuple \(\bar c\) in \(\mathcal U\). Let
\(D=\delta(\mathcal U,\bar c)\) be infinite and strongly minimal: in every
elementary extension of \((\mathcal U,\bar c)\), every subset of the
interpretation of \(D\) definable with arbitrary parameters is finite or
cofinite. All uses of \(\operatorname{acl}\) refer to the ambient structure
and its elementary extensions, not just to the language induced on \(D\).
For \(B\subseteq\mathcal U\), an element is algebraic over \(B\) if a
formula with a finite tuple of parameters from \(B\) defines a finite set
containing it. For \(X\subseteq D\), put
\[
 \operatorname{cl}_{\bar c}(X)
   \defeq D\cap\operatorname{acl}(\bar c\cup X).
\]
The notation \(\bar c\cup X\) means the set of coordinates of \(\bar c\)
together with \(X\). A set \(I\subseteq D\) is independent over a parameter
set \(B\supseteq\bar c\) when
\(a\notin\operatorname{acl}(B\cup(I\setminus\{a\}))\) for every \(a\in I\).
A basis of \(D\) over \(\bar c\) is a maximal such independent subset.
\end{hbconvention}

\begin{hblemma}{thm:p03-bl-uniform}{proved=annex:p3.n-a-006.2.uniform}{Uniform finite/cofinite bounds}
For every formula \(\theta(x,\bar y,\bar c)\), where \(x\) ranges over \(D\),
there is a natural number \(q\) such that, in every elementary extension and
for every \(\bar b\), either
\[
 \abs{\{x\in D:\theta(x,\bar b,\bar c)\}}\leq q
 \quad\text{or}\quad
 \abs{\{x\in D:\neg\theta(x,\bar b,\bar c)\}}\leq q.
\]
In particular, within this family infinitude, cofinitude, and having more
than \(q\) solutions are equivalent and are expressible by a formula in
\(\bar y\).
\end{hblemma}

The elementary-extension clause in strong minimality is essential here.
Compactness must exclude a parameter instance with infinitely many
solutions and infinitely many nonsolutions, not just instances visible in
one small model.

\begin{hbtheorem}{thm:p03-bl-pregeometry}{proved=annex:p3.n-a-006.2.pregeometry}{Algebraic closure gives a pregeometry}
The operation \(\operatorname{cl}_{\bar c}\) is extensive, monotone,
idempotent, and of finite character. For \(a,b\in D\) and \(X\subseteq D\),
it satisfies exchange:
\[
 a\in\operatorname{cl}_{\bar c}(X\cup\{b\})
       \setminus\operatorname{cl}_{\bar c}(X)
 \quad\Longrightarrow\quad
 b\in\operatorname{cl}_{\bar c}(X\cup\{a\}).
\]
The conclusion remains valid with any additional fixed parameter set in
the algebraic-closure base. No countability hypothesis is needed for this
theorem.
\end{hbtheorem}

The proof is a finite-fibre argument. Failure of exchange would make a
relation have cofinite fibres in one direction and uniformly bounded finite
fibres in the other. Intersecting finitely many cofinite fibres then
contradicts the finite bound. This is a proof of the difficult closure
axiom, not an assumption that the set is already a vector space.

\begin{hbtheorem}{thm:p03-bl-dimension}{proved=annex:p3.n-a-006.2.dimension}{Bases and dimension}
Assume the Axiom of Choice in the ambient set theory. Every independent
subset of \(D\) extends to a basis, each basis \(I\) spans
\(D=\operatorname{cl}_{\bar c}(I)\), and all bases have the same cardinality.
Write this cardinal as \(\dim_{\bar c}(D)\). More generally, an independent
set in the closure of a finite set of \(r\) elements has at most \(r\)
elements. If the language is countable and \(\bar c\) is finite, then
\[
 \abs{D}\leq\max(\aleph_0,\dim_{\bar c}(D)).
\]
Consequently an uncountable \(D\) has
\(\dim_{\bar c}(D)=\abs{D}\).
\end{hbtheorem}

\begin{hblemma}{thm:p03-bl-independent-types}{proved=annex:p3.n-a-006.2.independent-types}{The type of an independent tuple}
For every parameter set \(B\supseteq\bar c\), all elements of
\(D\setminus\operatorname{acl}(B)\) have the same complete type over \(B\).
Any two ordered independent tuples in \(D\) of the same finite length over
\(B\) have the same complete type over \(B\). Hence any bijection between
independent subsets over \(B\), together with the identity on \(B\), is an
elementary map on its domain.
\end{hblemma}

\subsection{The whole-universe case and its exact limit}

\begin{hbdefinition}{def:p03-bl-strongly-minimal-theory}{A strongly minimal theory}
A complete theory \(T\) with infinite models is strongly minimal when
\(x=x\) is strongly minimal in its models. Thus its entire universe, not
merely one definable subset, is a strongly minimal set. In this subsection
the base tuple is empty and every model \(M\) has a basis for
\(\operatorname{acl}\) on its whole universe.
\end{hbdefinition}

\begin{hbtheorem}{thm:p03-bl-whole-models}{proved=annex:p3.n-a-006.2.whole-models}{Classification and prime-over-basis in the whole-universe case}
Assume ambient choice, and let \(T\) be strongly minimal. Two models are isomorphic if and only if
their bases have the same cardinality; any bijection between bases extends
to an isomorphism. Every model \(M\) is prime and minimal over a basis
\(I\): any elementary map of \(I\) into a model of \(T\) extends to an
elementary embedding of \(M\), and an elementary submodel of \(M\)
containing \(I\) is \(M\) itself.

For a countable language, \(T\) is categorical in every uncountable
cardinal, and its countable-model spectrum is either one or \(\aleph_0\).
More precisely, if a finite-dimensional model exists, there is a least
finite dimension \(n_0\), and the countable models have exactly the
dimensions \(n_0,n_0+1,\ldots,\aleph_0\). If no finite-dimensional model
exists, every countable model has dimension \(\aleph_0\).
\end{hbtheorem}

\begin{proofboundary}{def:p03-bl-coordinatization-boundary}{Geometry is not yet general coordinatization}
The preceding construction is justified by
\(M=\operatorname{acl}(I)\) for the \emph{whole universe}. It does not
prove that, for an arbitrary uncountably categorical theory, a strongly
minimal \(D\) exists over parameters in a prime model or that every model
is prime and minimal over a basis of \(D\). Nor does prime-over-a-basis
imply that all elements are algebraic over that basis in this more general
setting. The existence and model-control assertions belong to the
Baldwin--Lachlan refinement and require their own argument
\cite{BaldwinLachlan1971}. The original article's general countable-model
argument and its fourth-section example are not replaced by the
whole-universe proof here.
\end{proofboundary}

\begin{hbexample}{def:p03-bl-extra-sort}{A strongly minimal subset that does not control the model}
Take the language with one unary predicate \(P\) and the complete theory
that both \(P\) and its complement are infinite pure sets. Its quantifier
elimination follows by extending finite partial bijections separately in
the two cells. The set \(P\) is strongly minimal: with finitely many
parameters, a definable subset of \(P\) is a finite or cofinite union of
equality cells. Nevertheless the full model also has the independent
cardinal invariant \(\abs{M\setminus P}\). In an uncountable cardinal
\(\kappa\), the models with cell sizes \((\kappa,\aleph_0)\) and
\((\kappa,\kappa)\) are not isomorphic, even though their \(P\)-dimensions
agree. The mere presence of a strongly minimal definable set therefore
does not supply the model-control step.
\end{hbexample}

\subsection{Fixed-characteristic algebraically closed fields}

\begin{hbtheorem}{thm:p03-bl-acf-geometry}{proved=annex:p3.n-a-006.2.acf}{Field algebraicity and the dimension illustration}
Fix \(p=0\) or a prime, and work in the ring language with
\(\mathsf{ACF}_p\). Using the quantifier-elimination theorem
\ref{thm:algebraically-closed-fields}, this is a complete strongly minimal theory.
If \(K\models\mathsf{ACF}_p\) and \(B\subseteq K\), then
\(\operatorname{acl}(B)\) is the set of elements of \(K\) algebraic over
the field generated by \(B\) and the prime field. A pregeometry basis is
therefore a transcendence basis. There is exactly one countable model for
each transcendence degree \(0,1,2,\ldots,\aleph_0\), and exactly one model
in each uncountable cardinal.
\end{hbtheorem}

For an infinite one-variable definable set in a strongly minimal universe,
Morley rank is \(1\) and Morley degree is \(1\): every definable proper
split has a finite side, and two cofinite subsets cannot be disjoint.
These invariants stay fixed while the dimension of a countable
algebraically closed field ranges through \(0,1,2,\ldots,\aleph_0\).
Thus neither rank nor degree is the dimension of the field's basis.
Characteristic must be fixed; fields of different characteristics are
not elementarily equivalent. This field-algebraicity illustration uses
the quantifier-elimination and algebraic-closure dependencies stated
above; it is distinct from the general Baldwin--Lachlan coordinatization
argument.

Morley's original categoricity argument and the strongly minimal
refinement are different source presentations. The former's rank and
saturation/transfer machinery is not renamed as a pregeometry proof, and
the local proofs above do not establish the refinement's general
existence theorem \cite{Morley1965,BaldwinLachlan1971}.


\section{NIP and o-minimality}
\label{sec:p03-nip-ominimal}

\begin{hbdefinition}{def:p03-nip}{Shattering, IP, and NIP}
Let \(T\) be complete and let \(\varphi(x;y)\) be a formula.  In a model
\(\mathcal M\entails T\), the parameter family
\(\{\varphi(\mathcal M;b):b\in(\dom\mathcal M)^{\abs y}\}\)
\emph{shatters} a finite set
\(A\subseteq(\dom\mathcal M)^{\abs x}\) when, for every \(A_0\subseteq A\),
some parameter tuple \(b\) satisfies
\(\varphi(\mathcal M;b)\cap A=A_0\).  The formula has the
\emph{independence property (IP) in \(T\)} when arbitrarily large finite sets are
shattered in some model of \(T\); it is \emph{NIP} otherwise.  The theory
\(T\) is NIP when every formula is NIP.  Formula-level shattering and
theory-level NIP are distinct standings; the latter is used in the dividing
line comparison below.
\index{independence property}
\index{NIP}
\index{shattering}
\end{hbdefinition}

\begin{hbdefinition}{def:p03-o-minimal}{O-minimality}
An expansion \(\mathcal M\) of a dense linear order without endpoints is
\emph{o-minimal} when every subset of \(\dom\mathcal M\) definable with
parameters is a finite union of points and intervals, with unbounded endpoints
allowed.  A complete theory of ordered structures is o-minimal when all its
models are o-minimal.  This condition controls unary definable sets; it is
strictly stronger than merely being NIP and does not imply stability.  The
distinction is needed for the DLO and real-closed-field rows below.
\index{o-minimality}
\end{hbdefinition}

\begin{hbdefinition}{def:p03-ominimal-cell}{Cells in an o-minimal structure}
Let \(\mathcal M\) be o-minimal.  In \(M\), a \emph{cell} is a point or an
open interval, with either endpoint allowed to be infinite.  Recursively, if
\(C\subseteq M^n\) is a cell, then cells over \(C\) are either graphs of
continuous definable functions \(C\to M\), or bands
\[
\{(x,y):x\in C\text{ and }f(x)<y<g(x)\},
\]
where each of \(f,g\) is either a continuous definable \(M\)-valued function
on \(C\) or one of the two constant boundaries \(-\infty,+\infty\), and
\(f<g\).  Cells are needed because arbitrary definable sets
in several variables need not be unions of boxes.  They are recursively
defined subsets of Cartesian powers of the ordered domain, not new sorts or
formulas, and they refine the unary point-and-interval definition of
o-minimality.  The refinement is used in the decomposition theorem.
\index{cell decomposition}
\end{hbdefinition}

\begin{hbtheorem}{thm:p03-ominimal-monotonicity}{cited=PillaySteinhorn1986}{O-minimal monotonicity theorem}
Let \(f:(a,b)\to M\) be definable in an o-minimal structure, with parameters
allowed.  There is a finite partition of \((a,b)\) into points and open
intervals such that on each interval \(f\) is continuous and is either
constant, strictly increasing, or strictly decreasing.
\end{hbtheorem}

\begin{hbtheorem}{thm:p03-ominimal-cell-decomposition}{cited=VanDenDries1998}{O-minimal cell decomposition}
For every finite family of definable subsets of \(M^n\), there is a finite
partition of \(M^n\) into cells such that each set in the family is a union of
cells.  If finitely many definable functions are also specified, the partition
can be refined so that each restriction to every cell contained in its domain
is continuous.
\end{hbtheorem}

For real closed fields in the ordered-ring language, definable sets are the
semialgebraic sets by quantifier elimination, so these theorems give finite
semialgebraic monotonicity and cell decompositions.  They explain the word
``tame'' in o-minimal geometry: the definition controls one-dimensional sets,
whereas cell decomposition propagates that control to all finite dimensions
\cite{PillaySteinhorn1986,VanDenDries1998,Marker2002}.

Stable theories are NIP; o-minimal theories are NIP; DLO and real closed
fields (RCF) are o-minimal and NIP but DLO is unstable.  The random graph has IP,
and a theory that is both simple and NIP is stable
\cite{Simon2015,VanDenDries1998,PillaySteinhorn1986,TentZiegler2012}.

\begingroup\small
\setlength{\tabcolsep}{4pt}
\begin{longtable}{@{}p{0.13\textwidth}p{0.33\textwidth}p{0.20\textwidth}p{0.23\textwidth}@{}}
\caption{Classification-theoretic dividing lines.}
\label{tab:p03-classification-overview}\\
\toprule
Dividing line & Exact defining condition & Canonical positive example &
Boundary and implication \\
\midrule
\endfirsthead
\caption[]{Classification-theoretic dividing lines (continued).}\\
\toprule
Dividing line & Exact defining condition & Canonical positive example &
Boundary and implication \\
\midrule
\endhead
\midrule\multicolumn{4}{r}{\emph{continued on next page}}\\
\endfoot
\endlastfoot
Stable & No formula has the order property in a complete theory. &
Algebraically closed fields \(\mathrm{ACF}_p\), for fixed characteristic \(p\). &
DLO is unstable; stable implies simple and NIP. \\
\addlinespace
Simple & Forking has local character with the cardinal bound of
Definition~\ref{def:p03-simple-theory}. &
The random graph; every stable theory. &
The random graph is unstable, so simple does not imply stable. \\
\addlinespace
NIP & No formula shatters arbitrarily large finite sets. &
DLO, RCF, and every stable theory. &
The random graph has IP; simple plus NIP implies stable. \\
\addlinespace
O-minimal & Every parameter-definable unary set is a finite union of points
and intervals. &
DLO and RCF. &
O-minimal implies NIP but not stability. \\
\bottomrule
\end{longtable}
\endgroup

The table is a survey map, not an implication-free glossary.  Its definitions
are supplied immediately above, and its example and implication claims are
sourced by \cite{Shelah1971Stability,TentZiegler2012,Kim1998,Simon2015,
VanDenDries1998,PillaySteinhorn1986}.
